\chapter{Classical Error-Control Codes}
\label{ch:classiccodes}

\begin{chapterintro}
Chapter~\ref{ch:infotheory} ended with a promise and a puzzle. Shannon proved that every
channel has a capacity and that, below it, codes exist whose error probability can be made as
small as we like; he did not say how to build them, and his random codes cannot be decoded in
any reasonable time. This chapter tells the story of the first fifty years of the answer: the
\emph{classical} codes, built from algebra and from finite-state machines, that put
error control into every modem, disk drive, compact disc, spacecraft and cellular phone long
before turbo and LDPC codes closed the last few decibels. We start with why coding pays and
what it costs, then develop linear block codes from the ground up: generator and parity-check
matrices, syndromes, Hamming distance and the bounds that limit what any code can do. Hamming
and Golay codes lead to the SECDED memory protection inside every server. A gentle
introduction to finite fields opens the door to cyclic codes, the CRC that guards almost every
packet ever sent, BCH codes and the Reed--Solomon codes of the CD, DVB, QR codes and RAID-6.
Interleaving and concatenation show how codes are combined to defeat bursts. The second half
turns to convolutional codes and the Viterbi algorithm, which carried Voyager's pictures,
802.11 and digital television, then to sequential decoding and Reed--Muller codes, and closes
with the protocols (ARQ and hybrid ARQ) that let a receiver simply ask again.
\end{chapterintro}

%======================================================================
\section{Why code?}
\label{sec:ch14:why}
%======================================================================

An uncoded binary link over a Gaussian channel is at the mercy of a single number. With
antipodal signalling the bit error probability is $\Q(\sqrt{2E_b/N_0})$
(Chapter~\ref{ch:baseband}); to reach $10^{-5}$ it needs $E_b/N_0=9.6$\,dB, and each further
factor of ten in reliability costs roughly another decibel. The only knobs are transmit
power, antenna gain and distance. Channel coding adds a new knob. By sending $n$ channel bits
for every $k$ information bits, with the extra $n-k$ bits computed from the information in a
cleverly structured way, the receiver can detect and correct errors that the channel has made.
The ratio $R=k/n$ is the \term{code rate}.

\subsection{The simplest codes}

The most obvious way to protect a bit is to repeat it. The \term{repetition code} of length
three sends $000$ for a zero and $111$ for a one, and the receiver takes a majority vote. On a
binary symmetric channel that flips each bit independently with probability $p$, the vote is
wrong only if two or three bits flip:
\begin{equation}
P_e = 3p^2(1-p)+p^3\approx 3p^2 .
\label{eq:ch14:rep3}
\end{equation}
With $p=10^{-2}$ the error rate falls to about $3\times10^{-4}$. But the code has rate $1/3$:
it triples the bandwidth and, at fixed transmitter power, gives each channel bit only a
third of the energy. Redo the comparison at equal energy per \emph{information} bit and the
repetition code turns out to be \emph{worse} than sending the bit once with three times the
energy. Repetition is the wrong kind of redundancy: it spends energy without buying distance
efficiently. The whole art of coding is to find redundancy that pays for itself.

The second-simplest code points the way. Append to every block of $k$ bits a single
\term{parity bit}, the modulo-2 sum of the block, so that every transmitted word has an even
number of ones. Any single error makes the parity odd and is \emph{detected}, at the modest
cost of one bit in $k+1$. Parity cannot say \emph{which} bit is wrong, but it was enough for
the paper-tape readers, the memory of early computers and the seven-bit ASCII terminals of the
1960s. Richard Hamming's insight, in 1950, was that several overlapping parity checks can
locate an error as well as detect it.

\subsection{Coding gain}

The benefit of a code is measured as \term{coding gain}: the reduction in $E_b/N_0$ needed to
reach a target error rate, compared with uncoded transmission of the same information at the
same modulation. The comparison must be made at equal energy per \emph{information} bit, because
the code spends energy on redundant bits: each coded bit carries energy $E_c=RE_b$, so the
channel error rate before decoding is \emph{higher} than for an uncoded system. A code earns its
keep only if its error correction more than repays that loss.

Figure~\ref{fig:ch14_coding_gain} shows the landscape of the classical codes on the
additive white Gaussian noise (AWGN) channel with BPSK. At a bit error rate of $10^{-5}$, the
Hamming (7,4) code with hard decisions gains well under a decibel; the Golay code and a long BCH
code gain about 2 and 3.5\,dB; the rate-1/2 convolutional code with constraint length 7 and soft
Viterbi decoding gains about 5.5\,dB; and the concatenation of that code with a Reed--Solomon
outer code, the combination that carried pictures back from Uranus and Neptune, gains about 7\,dB
and reaches $10^{-5}$ at roughly 2.3\,dB with ideal interleaving (about 2.5\,dB is the figure
usually quoted for the deployed system). The Shannon limit for rate~$1/2$ binary signalling,
0.19\,dB (Chapter~\ref{ch:infotheory}), sits a further 2\,dB or so to the left. That remaining gap
is the subject of Chapter~\ref{ch:moderncodes}.

\mdcfig[0.88]{ch14_coding_gain}{The classical codes on the AWGN channel with BPSK. Block codes
with hard-decision decoding (dashed) are computed from the standard approximation
of~\eqref{eq:ch14:pbblock}; the convolutional and concatenated curves are Monte Carlo
simulations with the union bound (dotted) continuing the convolutional curve to low error rates.
The concatenated curve assumes ideal interleaving between the codes. Coding gain is read
horizontally at the target error rate.}

It helps to separate two numbers. The \term{asymptotic coding gain}, the gain as
$E_b/N_0\to\infty$, depends only on the rate and the minimum distance $d$ of the code (defined in
Section~\ref{sec:ch14:block}). With soft-decision maximum-likelihood decoding it is
\begin{equation}
G_{\infty}^{\text{soft}}=10\log_{10}(R\,d)\ \text{dB},\qquad
G_{\infty}^{\text{hard}}\approx10\log_{10}\!\big(R\,(t+1)\big)\ \text{dB},
\label{eq:ch14:acg}
\end{equation}
where $t=\lfloor(d-1)/2\rfloor$ is the number of errors a hard-decision decoder can correct. The
\emph{real} coding gain at a practical error rate is smaller, because codes have many nearest
neighbours and the asymptote is approached slowly. For the $K=7$ convolutional code, $R=1/2$ and
$d=10$ give an asymptotic gain of 7\,dB, of which about 5.5\,dB is realised at $10^{-5}$ and
almost 6\,dB at $10^{-7}$ (Section~\ref{sec:ch14:viterbi}).

\subsection{What coding costs}

Coding gain is not free. It is paid for in four currencies, and a system designer chooses a
code by deciding which of them is cheapest.
\begin{description}[style=nextline,leftmargin=1.2em,font=\sffamily\bfseries\small\color{navy}]
\item[Bandwidth or throughput.] A rate-$R$ code sends $1/R$ channel symbols per information bit.
At a fixed symbol rate the throughput falls by $R$; at a fixed throughput the bandwidth grows by
$1/R$. In \emph{power-limited} systems (deep space, satellites, low-power IoT) bandwidth is
cheap and low-rate codes are ideal; in \emph{bandwidth-limited} systems (cable, microwave
backhaul, dense cellular) the code must be combined with a larger constellation, the subject of
coded modulation (Chapters~\ref{ch:modulation} and~\ref{ch:moderncodes}).
\item[Latency.] A block decoder cannot finish until the whole block has arrived; a Viterbi
decoder must look ahead several constraint lengths; an interleaver adds its full span. Voice and
control loops tolerate milliseconds, not seconds.
\item[Complexity and power.] Decoding costs gates, memory and energy. In the 1970s a
Viterbi decoder for $K=7$ filled a rack; today it occupies a fraction of a square millimetre, but
at 400\,Gb/s the decoder's power budget still limits what code a link can afford.
\item[Error behaviour.] Every decoder eventually fails. A block decoder may detect the failure
or may \emph{miscorrect}, confidently delivering a wrong codeword; a Viterbi decoder fails in
bursts. The code must be matched to what the next layer can tolerate.
\end{description}

\subsection{Three strategies: FEC, ARQ and HARQ}

There are two fundamentally different ways to use redundancy. \term{Forward error correction}
(FEC) sends enough redundancy for the receiver to correct errors on its own, with no return
channel. \term{Automatic repeat request} (ARQ) sends only enough redundancy to \emph{detect}
errors, usually a cyclic redundancy check, and asks the transmitter to try again when a check
fails. ARQ is extraordinarily efficient when errors are rare, because detection is cheap: 32
check bits protect a 12\,000-bit Ethernet frame against virtually every error pattern. But ARQ
needs a return channel and adds a round trip of delay for every failure, which rules it out for
broadcast, for deep space (a round trip to Neptune takes about eight hours) and for storage,
where the ``transmitter'' may have been a disk head that wrote the data years ago.

Most modern wireless links combine the two in \term{hybrid ARQ} (HARQ): an FEC code corrects
most errors, a CRC detects the residue, and retransmissions carry either a copy of the codeword
(Chase combining) or fresh parity bits (incremental redundancy) that the receiver combines with
what it already has. Section~\ref{sec:ch14:arq} develops these protocols; Chapter~\ref{ch:cellular}
shows them at work in LTE and 5G.

\begin{keyidea}[Coding buys distance, and distance buys reliability]
A code is a set of points chosen from a much larger space so that they are far apart. The
channel pushes the transmitted point around; the decoder picks the nearest codeword. Everything
in this chapter follows from three questions: how far apart can the codewords be for a given rate
(the bounds), how can we build codes that achieve large distance with enough structure to decode
(the algebra and the trellis), and how can we find the nearest codeword quickly (the decoding
algorithms)?
\end{keyidea}

\begin{history}[Hamming's weekends]
In 1947 Richard Hamming was a young mathematician at Bell Telephone Laboratories with access to
the Bell Labs Model~V relay computer, but only at weekends, when it ran unattended jobs from a
queue. The machine checked its own arithmetic with two-out-of-five codes and, when a check failed,
simply abandoned the job and moved to the next. Returning on two consecutive Mondays to find his
work discarded, Hamming later recalled asking himself, with some irritation, why a machine that
could \emph{detect} an error could not \emph{locate} it and fix it. His answer, published in the
\emph{Bell System Technical Journal} in April 1950 (Bell Labs held up publication until a
patent was filed), introduced the single-error-correcting codes that bear his name, the
geometric notion of distance between binary words and the idea of a code as a packing of
spheres. Marcel Golay at the Signal Corps laboratories had seen Hamming's (7,4) code in Shannon's
1948 paper and in 1949 published a one-page note generalising it and describing the remarkable
(23,12) code that now bears his name.
\end{history}

\subsection{Seventy years in one table}

Table~\ref{tab:ch14:timeline} lists the milestones of the classical era. Two threads run
through it. The \emph{algebraic} thread (Hamming, Golay, Reed--Muller, BCH, Reed--Solomon)
builds block codes with large guaranteed distance and decodes them by solving equations over
finite fields. The \emph{probabilistic} thread (Elias, Wozencraft, Fano, Viterbi) builds
convolutional codes, generated by shift registers, and decodes them by searching a tree or
trellis for the most likely path. Forney's concatenation joined the two threads in 1966, and
their descendants---turbo and LDPC codes---closed the gap to capacity in the 1990s.

\begin{table}[htbp]\centering\small
\caption{Milestones of classical coding.}
\label{tab:ch14:timeline}
\begin{tabularx}{\linewidth}{@{}l X@{}}
\toprule
Year & Milestone\\
\midrule
1948 & Shannon's channel coding theorem; the (7,4) Hamming code appears as an example in his paper\\
1949 & Golay's note on the perfect binary (23,12) and ternary (11,6) codes\\
1950 & Hamming, ``Error detecting and error correcting codes''\\
1954 & Reed--Muller codes (Muller: construction; Reed: majority-logic decoding); Elias's product codes\\
1955 & Elias introduces convolutional codes\\
1957 & Prange's cyclic codes; Wozencraft's sequential decoding\\
1959--60 & BCH codes (Hocquenghem; Bose and Ray-Chaudhuri); Reed--Solomon codes; Peterson's decoder\\
1963 & Fano's sequential-decoding algorithm; Gallager's LDPC codes (ahead of their time)\\
1966 & Forney's concatenated codes and generalised minimum-distance decoding\\
1967 & Viterbi's algorithm\\
1968--69 & Berlekamp's decoding algorithm; Massey's shift-register-synthesis interpretation\\
1969--71 & Mariner 6, 7 and 9 send Mars images with the (32,6) Reed--Muller code\\
1972--73 & Pioneer 10 and 11 use a $K=32$ convolutional code with sequential decoding\\
1977 & Voyager launches with a $K=7$, rate-1/2 convolutional code; RS(255,223) concatenation added for the outer planets\\
1980--82 & Compact disc with cross-interleaved Reed--Solomon coding (CIRC)\\
1987 & CCSDS standardises RS(255,223) concatenated with the $K=7$ code\\
1993--95 & Turbo codes; LDPC codes rediscovered (Chapter~\ref{ch:moderncodes}); DVB-S uses RS(204,188) + $K=7$\\
1994 & QR code (Denso Wave) with Reed--Solomon error correction\\
\bottomrule
\end{tabularx}
\end{table}

%======================================================================
\section{Linear block codes}
\label{sec:ch14:block}
%======================================================================

An $(n,k)$ \term{block code} maps each block of $k$ information bits to a \emph{codeword} of
$n$ bits. Of the $2^n$ possible $n$-bit words, only $2^k$ are codewords. The receiver sees an
$n$-bit word that may have been corrupted and must decide which codeword was sent. With no
further structure, encoding would require a table of $2^k$ codewords and decoding a search over
all of them---impossible for the $k=223\times8=1784$-bit blocks of a Reed--Solomon code. The
structure that makes coding practical is linearity.

\subsection{Linearity, generator and parity-check matrices}

A code is \term{linear} if the modulo-2 sum of any two codewords is a codeword. The codewords
then form a $k$-dimensional subspace of the vector space $\{0,1\}^n$ over the binary field
GF(2), where addition is exclusive-OR and multiplication is AND. Any basis of the subspace, $k$
linearly independent codewords written as the rows of a $k\times n$ matrix $\vect{G}$, generates
the code:
\begin{equation}
\vect{c}=\vect{u}\,\vect{G},\qquad \vect{u}\in\{0,1\}^k .
\label{eq:ch14:enc}
\end{equation}
$\vect{G}$ is the \term{generator matrix}. Encoding is now a matrix multiplication, and the code
is described by $kn$ bits instead of $n2^k$. Row operations on $\vect{G}$ change the mapping
from messages to codewords but not the code itself, so every linear code has a generator matrix
in \term{systematic form},
\begin{equation}
\vect{G}=\big[\,\vect{I}_k \mid \vect{P}\,\big],
\end{equation}
possibly after permuting columns. A systematic codeword is the message followed by $n-k$
\term{parity bits} $\vect{u}\vect{P}$; the receiver can read the message directly, which is
convenient and is what nearly every practical code does.

The same subspace can be described by the constraints it satisfies. A $(n-k)\times n$
\term{parity-check matrix} $\vect{H}$ has the codewords as its null space:
\begin{equation}
\vect{c}\,\vect{H}^{\mathsf T}=\vect{0}\quad\Longleftrightarrow\quad \vect{c}\ \text{is a codeword}.
\label{eq:ch14:parity}
\end{equation}
For a systematic generator, $\vect{H}=[\,\vect{P}^{\mathsf T}\mid\vect{I}_{n-k}\,]$, and
$\vect{G}\vect{H}^{\mathsf T}=\vect{P}+\vect{P}=\vect{0}$ confirms that every row of $\vect{G}$
satisfies every check. Each row of $\vect{H}$ is one parity equation: a set of code bits whose
sum must be zero. The code defined by $\vect{H}$ as a generator matrix is the \term{dual code},
of dimension $n-k$.

For the Hamming (7,4) code used throughout this chapter (and in \lab{commlib/coding.py}),
\begin{equation}
\vect{G}=\left[\begin{array}{cccc|ccc}
1&0&0&0&1&1&0\\ 0&1&0&0&1&0&1\\ 0&0&1&0&0&1&1\\ 0&0&0&1&1&1&1
\end{array}\right],\qquad
\vect{H}=\left[\begin{array}{cccc|ccc}
1&1&0&1&1&0&0\\ 1&0&1&1&0&1&0\\ 0&1&1&1&0&0&1
\end{array}\right].
\label{eq:ch14:h74}
\end{equation}
The three parity bits are $p_1=u_1\oplus u_2\oplus u_4$, $p_2=u_1\oplus u_3\oplus u_4$ and
$p_3=u_2\oplus u_3\oplus u_4$: three overlapping checks, each covering a different subset of the
data. Hamming's own drawing was a Venn diagram of three circles, with the four data bits in the
overlaps and each parity bit completing the even parity of its circle.

\subsection{Distance and weight}

The \term{Hamming distance} $d(\vect{x},\vect{y})$ between two words is the number of positions
in which they differ; the \term{Hamming weight} $w(\vect{x})$ is the number of ones. Distance is
a true metric (it satisfies the triangle inequality), and on the binary symmetric channel with
$p<1/2$ the most likely codeword is the one at the smallest Hamming distance from the received
word. The \term{minimum distance} $d_{\min}$ of a code is the smallest distance between two
distinct codewords, and it is the single most important number describing a code. We write an
$(n,k)$ code with minimum distance $d$ as an $(n,k,d)$ code.

For a linear code, $d(\vect{c}_1,\vect{c}_2)=w(\vect{c}_1\oplus\vect{c}_2)$, and
$\vect{c}_1\oplus\vect{c}_2$ is itself a codeword. So
\begin{equation}
d_{\min}=\min_{\vect{c}\ne\vect{0}} w(\vect{c}),
\end{equation}
the minimum weight of a nonzero codeword. There is another useful characterisation: since
$\vect{c}\vect{H}^{\mathsf T}=\vect{0}$ says that the columns of $\vect{H}$ selected by the ones
of $\vect{c}$ sum to zero, $d_{\min}$ is the smallest number of columns of $\vect{H}$ that are
linearly dependent. If no two columns of $\vect{H}$ are equal and none is zero, any single column
is nonzero and no two sum to zero, so $d_{\min}\ge3$. The seven columns of $\vect{H}$
in~\eqref{eq:ch14:h74} are exactly the seven nonzero 3-bit vectors, so the (7,4) code has
$d_{\min}=3$.

\subsection{What distance buys}

Draw a sphere of radius $t$ (in Hamming distance) around each codeword. If $d_{\min}\ge2t+1$ the
spheres do not overlap, so a received word with at most $t$ errors lies inside the sphere of the
transmitted codeword and is decoded correctly. If $d_{\min}\ge s+1$, no pattern of $s$ or fewer
errors can turn one codeword into another, so every such pattern is detected. Combining the two,
and adding \term{erasures}---positions the receiver knows are unreliable but whose values are
unknown---a code with minimum distance $d$ can
\begin{equation}
\text{correct } t \text{ errors and } e \text{ erasures whenever } 2t+e<d,
\label{eq:ch14:errera}
\end{equation}
or correct $t$ errors and simultaneously detect up to $s\ge t$ errors whenever $t+s<d$.
An erasure costs half as much as an error because the decoder knows where it is and only has to
find its value. Figure~\ref{fig:ch14_cube} shows the two smallest interesting codes on the
cube of 3-bit words.

\mdcfig[0.9]{ch14_cube}{Codes as sets of points on the cube $\{0,1\}^3$, where Hamming
distance is the number of edges between vertices. Left: the repetition code $\{000,111\}$
has $d_{\min}=3$; every word is within distance one of exactly one codeword, so single errors are
corrected. Right: the even-parity code has $d_{\min}=2$; every single error lands on an odd-weight
word that is not a codeword and is detected, but it is equidistant from three codewords, so it
cannot be corrected.}

\subsection{Syndrome decoding}

Let the received word be $\vect{r}=\vect{c}\oplus\vect{e}$, where $\vect{e}$ is the error
pattern. Multiply by $\vect{H}^{\mathsf T}$:
\begin{equation}
\vect{s}=\vect{r}\,\vect{H}^{\mathsf T}=\vect{c}\,\vect{H}^{\mathsf T}\oplus\vect{e}\,\vect{H}^{\mathsf T}
=\vect{e}\,\vect{H}^{\mathsf T}.
\label{eq:ch14:syndrome}
\end{equation}
The $(n-k)$-bit \term{syndrome} $\vect{s}$ depends only on the error pattern, not on the codeword
that was sent. It is zero when $\vect{r}$ is a codeword (no detectable error) and otherwise tells
the decoder which parity checks failed. If a single error occurs in position $j$, the syndrome
is the $j$-th column of $\vect{H}$. For a Hamming code, whose columns are all distinct, the
syndrome therefore points directly at the erroneous bit.

\begin{worked}[Decoding a Hamming (7,4) word by syndrome]
Encode $\vect{u}=1011$ with~\eqref{eq:ch14:h74}. The parity bits are
$p_1=1\oplus0\oplus1=0$, $p_2=1\oplus1\oplus1=1$, $p_3=0\oplus1\oplus1=0$, so
$\vect{c}=1011\,010$. Suppose the channel flips the sixth bit, giving $\vect{r}=1011\,000$.

The syndrome is the sum of the columns of $\vect{H}$ at the positions of the ones in
$\vect{r}$ (positions 1, 3 and 4):
\[
\vect{s}=\begin{bmatrix}1\\1\\0\end{bmatrix}\oplus\begin{bmatrix}0\\1\\1\end{bmatrix}\oplus
\begin{bmatrix}1\\1\\1\end{bmatrix}=\begin{bmatrix}0\\1\\0\end{bmatrix}.
\]
Checks 1 and 3 pass and check 2 fails. The vector $010$ is the sixth column of $\vect{H}$, so the
decoder flips bit 6 and recovers $1011\,010$ and the message $1011$. Had the error hit a
data bit, say bit 4, the syndrome would have been $111$ (column 4), because $u_4$ appears in all
three checks. With \emph{two} errors, for example in bits 1 and 2, the syndrome is
$110\oplus101=011$, the column for bit 3: the decoder confidently flips a third, correct bit
and outputs a codeword at distance three from the one sent. Hamming codes always correct
single errors and always miscorrect double errors.
\end{worked}

In general, many error patterns share each syndrome: exactly $2^k$ of them, one for each
codeword, since $\vect{e}$ and $\vect{e}\oplus\vect{c}$ give the same syndrome. These sets are
the \term{cosets} of the code. The maximum-likelihood decoder on a binary symmetric channel
chooses, for each syndrome, the lowest-weight pattern in the coset---the \term{coset leader}---and
adds it to $\vect{r}$. Arranging all $2^n$ words in a table with one coset per row, leaders in the
first column and the codewords along the top, gives the \term{standard array}. It is a
beautiful teaching device and a hopeless implementation for all but tiny codes: a (63,45) BCH code
has $2^{18}$ syndromes, which is manageable as a lookup table, but a (255,223) Reed--Solomon code
has $2^{256}$. Practical decoders compute the coset leader from the syndrome algebraically,
which is exactly what the algebraic structure of the later sections is for.

\subsection{Weight distribution and the error probabilities of a block code}

The number $A_w$ of codewords of weight $w$ is the code's \term{weight distribution}. It
determines the probability of every kind of decoding outcome. An error pattern goes
\emph{undetected} exactly when it is itself a nonzero codeword, so on a BSC with crossover
probability $p$
\begin{equation}
P_{\text{ud}}=\sum_{w=d_{\min}}^{n}A_w\,p^w(1-p)^{n-w}.
\label{eq:ch14:pud}
\end{equation}
For the Hamming (7,4) code, $A_3=7$, $A_4=7$, $A_7=1$, and $P_{\text{ud}}\approx7p^3$. A useful
identity due to MacWilliams relates the weight distribution of a code to that of its dual; it
makes~\eqref{eq:ch14:pud} computable for CRCs, whose duals are small (Section~\ref{sec:ch14:crc}).

For a code that corrects all patterns of up to $t$ errors and nothing more (a
\emph{bounded-distance} decoder), a decoded block is wrong whenever more than $t$ errors occur.
A standard approximation to the resulting bit error rate assumes that a failed block with $i$
channel errors leaves about $i+t$ bit errors in the $n$-bit block:
\begin{equation}
P_b\approx\sum_{i=t+1}^{n}\frac{i+t}{n}\binom{n}{i}p^i(1-p)^{n-i}.
\label{eq:ch14:pbblock}
\end{equation}
The dashed curves of Figure~\ref{fig:ch14_coding_gain} use this formula with
$p=\Q\big(\sqrt{2RE_b/N_0}\big)$.

\subsection{Making new codes from old}

Practical systems rarely use a code at its natural length. A handful of simple modifications
adjust $n$ and $k$ to fit a frame, summarised in Table~\ref{tab:ch14:modify}.

\begin{table}[htbp]\centering\small
\caption{Modifying an $(n,k,d)$ code.}
\label{tab:ch14:modify}
\begin{tabularx}{\linewidth}{@{}l l X@{}}
\toprule
Operation & New parameters & How and where\\
\midrule
Shortening & $(n-s,\,k-s,\,\ge d)$ & Fix $s$ message bits to zero and do not send them. RS(204,188) in DVB is RS(255,239) shortened by 51 bytes; every CRC is a shortened cyclic code.\\
Puncturing & $(n-s,\,k,\,\ge d-s)$ & Delete $s$ parity bits. Raises the rate; the decoder treats them as erasures. Used heavily with convolutional and turbo codes.\\
Extending & $(n+1,\,k,\,d+1)$ for odd $d$ & Add an overall parity bit. Hamming (7,4,3) becomes (8,4,4); Golay (23,12,7) becomes (24,12,8).\\
Expurgating & $(n,\,k-1,\,\ge d)$ & Keep only the even-weight codewords.\\
Lengthening & $(n+1,\,k+1,\,\le d)$ & Add a message bit (the inverse of shortening).\\
\bottomrule
\end{tabularx}
\end{table}

%======================================================================
\section{The limits of block codes}
\label{sec:ch14:bounds}
%======================================================================

How large can $d_{\min}$ be for given $n$ and $k$? The question has occupied coding theorists
since 1950 and is still open in general. Four classical bounds frame the answer, and they are
worth knowing because they tell an engineer immediately whether a proposed code is good, bad or
impossible.

\subsection{The Hamming (sphere-packing) bound}

A code that corrects $t$ errors surrounds each of its $2^k$ codewords with a disjoint sphere
containing $V(n,t)=\sum_{i=0}^{t}\binom ni$ words. The spheres must fit into the space of $2^n$
words, so
\begin{equation}
2^k\sum_{i=0}^{t}\binom ni\le2^n
\qquad\Longleftrightarrow\qquad
n-k\ \ge\ \log_2\sum_{i=0}^{t}\binom ni .
\label{eq:ch14:hamming}
\end{equation}
Each of the $2^{n-k}$ syndromes can name at most one correctable error pattern. A code that meets
the bound with equality, whose spheres tile the whole space with no word left over, is
\term{perfect}. Perfect codes are astonishingly rare. Apart from trivial cases (a single
codeword, the whole space, odd-length binary repetition codes), the only perfect binary codes are
the Hamming codes ($t=1$) and the (23,12) Golay code ($t=3$), a fact proved by Tietäväinen and by
van Lint in the early 1970s. For the Golay code, $1+23+253+1771=2048=2^{11}$ exactly: the
spheres of radius three around its 4096 codewords fill the 8\,388\,608 words of length 23 with
nothing to spare.

\subsection{The Singleton bound and MDS codes}

Delete any $d-1$ coordinates from every codeword. The shortened words are still distinct (two
codewords differ in at least $d$ places, so at least one difference survives), and there are only
$2^{n-d+1}$ possible words of the remaining length. Hence $2^k\le2^{n-d+1}$, or
\begin{equation}
d\le n-k+1.
\label{eq:ch14:singleton}
\end{equation}
The argument holds for codes over any alphabet. Codes meeting the \term{Singleton bound} with
equality are called \term{maximum-distance separable} (MDS): every $k$ symbols of a codeword
determine the rest, so the codeword can be recovered from \emph{any} $k$ of its $n$ symbols. Binary
MDS codes are trivial (repetition and single-parity-check codes), but over large alphabets they
abound: the Reed--Solomon codes of Section~\ref{sec:ch14:rs} are MDS, which is why they dominate
erasure coding in storage.

\subsection{The Gilbert--Varshamov bound}

The previous bounds say what cannot be done. The \term{Gilbert--Varshamov bound} says what
\emph{can}: a linear $(n,k)$ code with minimum distance at least $d$ exists whenever
\begin{equation}
2^{n-k}>\sum_{i=0}^{d-2}\binom{n-1}{i}.
\label{eq:ch14:gv}
\end{equation}
The proof builds $\vect{H}$ one column at a time, choosing each new column so that it is not the
sum of $d-2$ or fewer of the columns already chosen; the counting shows a choice always exists
while~\eqref{eq:ch14:gv} holds. It is a random-coding style argument in miniature, and
like Shannon's theorem it guarantees existence without construction. Remarkably, for long binary
codes no explicit family is known that beats it asymptotically.

\subsection{The Plotkin bound, and the asymptotic picture}

When $d$ is a large fraction of $n$, a simple averaging argument gives a tighter limit. The sum of
all pairwise distances among $M$ codewords is at most $nM^2/4$ (each coordinate contributes at
most $M^2/4$ disagreeing ordered pairs), and at least $M(M-1)d$. Hence, if $2d>n$,
\begin{equation}
M\le\frac{2d}{2d-n},
\label{eq:ch14:plotkin}
\end{equation}
and asymptotically no family with relative distance $\delta=d/n\ge1/2$ can have positive rate.

Letting $n\to\infty$ with $\delta$ fixed turns the bounds into curves of rate against relative
distance (Figure~\ref{fig:ch14_bounds}, left). Using $\log_2\binom{n}{\delta n}\approx
nH_2(\delta)$, where $H_2$ is the binary entropy function of Chapter~\ref{ch:infotheory}, the
Hamming bound becomes $R\le1-H_2(\delta/2)$ and the Gilbert--Varshamov bound $R\ge1-H_2(\delta)$.
The best known upper bound, from the linear-programming method of McEliece, Rodemich, Rumsey and
Welch (MRRW, 1977), lies between them, and the true answer lies somewhere in the gap. The right
panel shows the same bounds for $n=63$, with the BCH codes of Section~\ref{sec:ch14:bch}: for small
$d$ they sit right against the Hamming bound; for large $d$ they fall towards the
Gilbert--Varshamov curve.

\mdcfig{ch14_bounds}{Bounds on block codes. Left: asymptotic bounds on rate versus relative
distance; codes meeting the Gilbert--Varshamov bound are guaranteed to exist (shaded), and no code
can exceed the lowest upper bound (Hamming for small $\delta$, MRRW above about $0.3$). Right:
the dimension achievable at length 63 against minimum distance; the dots are the primitive BCH
codes.}

\begin{keyidea}[Minimum distance is not the whole story]
The bounds concern minimum distance, the right figure of merit for hard-decision
bounded-distance decoding at very low error rates. At the error rates where real systems operate,
what matters is the whole weight distribution and how close the decoder comes to maximum
likelihood. Shannon's random codes have mediocre minimum distance but superb performance, and
turbo and LDPC codes (Chapter~\ref{ch:moderncodes}) have relatively small minimum distances yet
come within a fraction of a decibel of capacity. The classical codes were designed for distance;
the modern codes are designed for decoding.
\end{keyidea}

%======================================================================
\section{Hamming, SECDED and Golay codes}
\label{sec:ch14:hamming}
%======================================================================

\subsection{The Hamming family}

The construction of the (7,4) code generalises at once. For any $m\ge2$, take as the columns of
$\vect{H}$ all $2^m-1$ nonzero binary vectors of length $m$. No two columns are equal, so
$d_{\min}\ge3$; three columns such as $001$, $010$ and $011$ sum to zero, so $d_{\min}=3$ exactly. The
result is the \term{Hamming code}
\begin{equation}
(n,k,d)=(2^m-1,\;2^m-1-m,\;3),
\end{equation}
giving (3,1), (7,4), (15,11), (31,26), (63,57), (127,120) and so on. Each is perfect: the
$2^m$ syndromes are exactly the zero syndrome plus one for each of the $n$ single-error
positions. If the columns are arranged in natural binary order ($001,010,011,\dots$), the
syndrome of a single error \emph{is} the binary index of the erroneous bit, which is how Hamming
described it. The rate $1-m/(2^m-1)$ approaches one as $m$ grows, but the correcting power stays at
one error per block, so long Hamming codes are useful only where errors are very rare.

\subsection{SECDED: extended Hamming codes in computer memory}

Adding an overall parity bit to a Hamming code gives the \term{extended Hamming code}
$(2^m,2^m-1-m,4)$. With $d_{\min}=4$ it can correct one error and at the same time detect any two:
a \term{SECDED} (single-error-correcting, double-error-detecting) code. The decoder looks at the
Hamming syndrome $\vect{s}$ and the overall parity $q$:
\begin{itemize}
\item $\vect{s}=\vect{0}$, $q=0$: no error;
\item $q=1$: an odd number of errors, assumed to be one; correct the bit named by $\vect{s}$ (or the
parity bit itself if $\vect{s}=\vect{0}$);
\item $\vect{s}\ne\vect{0}$, $q=0$: an even number of errors, at least two; report an
uncorrectable error.
\end{itemize}
SECDED is the workhorse of memory protection. A standard ECC memory module is 72 bits wide,
carrying 64 data bits and 8 check bits: a (72,64) shortened extended Hamming code. Hsiao's 1970
variant chooses the columns of $\vect{H}$ to have odd weight and the rows to have nearly equal
weight, which makes the syndrome logic shallower and faster and turns the ``even number of
errors'' test into a simple check on the syndrome's weight. The same code structure protects
processor caches, register files and on-chip buses, and DDR5 memory chips carry an internal
single-error-correcting (136,128) code inside the die, in addition to any module-level ECC.

\begin{inpractice}[Why your server needs ECC]
Memory bits flip. The causes include alpha particles from trace radioactivity in packaging,
neutrons produced by cosmic rays in the upper atmosphere (the rate at the altitude of Denver is
several times that at sea level), and, increasingly, cell-to-cell interference in densely packed
DRAM. A large field study of Google's server fleet by Schroeder, Pinheiro and Weber (2009) found
DRAM error rates orders of magnitude above earlier laboratory estimates, with roughly a third of
machines seeing at least one correctable error in a year. The ``Rowhammer'' disturbance errors
reported in 2014 showed that repeatedly accessing one DRAM row can flip bits in its neighbours.
SECDED turns almost all of these into logged, corrected events. Data centres go further with
\emph{Chipkill}-class codes, symbol-based codes (often Reed--Solomon-like) in which each symbol is
the group of bits from one DRAM chip, so that the complete failure of a chip is a single
correctable symbol error.
\end{inpractice}

\begin{pitfall}[SECDED silently miscorrects some triple errors]
The extended Hamming code has $d_{\min}=4$, so three errors can move the received word to within
distance one of a \emph{different} codeword. The decoder then sees odd parity, a nonzero
syndrome, and ``corrects'' a fourth bit, delivering wrong data with no warning. Multi-bit upsets
from a single particle strike are common in dense memories, which is why physically adjacent
bits are placed in different codewords (bit interleaving across words) and why the probability of
a triple error in one word, not the probability of a single error, sets the silent-data-corruption
rate of a memory system.
\end{pitfall}

\subsection{The Golay codes}

The binary \term{Golay code} is the $(23,12,7)$ perfect code that corrects any three errors in 23
bits. It is a cyclic code (Section~\ref{sec:ch14:cyclic}) with generator polynomial
$g(x)=x^{11}+x^{10}+x^6+x^5+x^4+x^2+1$, and its extension by an overall parity bit, the
$(24,12,8)$ \emph{extended Golay code}, is one of the most symmetric objects in mathematics: its
automorphism group is the Mathieu group $M_{24}$, and it leads, through the Leech lattice, to the
densest known sphere packing in 24 dimensions and to the classification of finite simple groups.
Engineers value it for humbler reasons. Rate $1/2$, triple-error correction and a small, fixed
block size make it ideal for protecting short headers and control words. The Voyager spacecraft
used the extended Golay code for their colour images of Jupiter and Saturn; the (24,12) code
protects the words of the automatic link establishment protocol of HF radio
(MIL-STD-188-141). Its hard-decision
coding gain at $10^{-5}$ is about 2.1\,dB, compared with 0.4\,dB for Hamming (7,4)
(Figure~\ref{fig:ch14_coding_gain}).

%======================================================================
\section{Hard and soft decisions}
\label{sec:ch14:softhard}
%======================================================================

So far the receiver has sliced each channel symbol to a bit before decoding. That is
\term{hard-decision decoding}: the channel as seen by the decoder is a binary symmetric channel
with crossover probability $p=\Q(\sqrt{2RE_b/N_0})$. But the demodulator had more to offer. A
received sample of $+0.05$ and one of $+1.9$ both slice to a one, yet the first is almost a coin
toss and the second is nearly certain. A \term{soft-decision} decoder uses the unquantised (or
finely quantised) samples, or equivalently the log-likelihood ratios
\begin{equation}
L_i=\ln\frac{P(c_i=0\mid y_i)}{P(c_i=1\mid y_i)}=\frac{2y_i}{\sigma^2}
\label{eq:ch14:llr}
\end{equation}
for BPSK in Gaussian noise of variance $\sigma^2$ (the convention of \lab{commlib}, bit 0 mapped to
$+1$). The maximum-likelihood soft decoder chooses the codeword $\vect{c}$ whose BPSK image
$\vect{x}=1-2\vect{c}$ is closest in \emph{Euclidean} distance to $\vect{y}$, or equivalently
has the largest correlation $\sum_i x_iy_i$.

Euclidean distance changes the arithmetic of distance. Two codewords at Hamming distance $d$ have
BPSK images at squared Euclidean distance $4dE_c=4dRE_b$, so the pairwise error probability is
$\Q\big(\sqrt{2dRE_b/N_0}\big)$: soft decoding exploits the \emph{full} distance $d$, whereas a
hard-decision decoder is guaranteed only $t+1\approx d/2$ errors' worth. This is the origin of the
3\,dB difference between the asymptotic gains in~\eqref{eq:ch14:acg}.

At practical error rates the difference is smaller, about 2\,dB, as Figure~\ref{fig:ch14_hard_soft}
shows for two codes: about 1.4\,dB for the short Hamming code and 2.3\,dB for the convolutional
code at $10^{-5}$. Information theory gives the same number from a different direction. At low
signal-to-noise ratio, slicing BPSK to a binary symmetric channel reduces capacity by a factor
$2/\pi$, a loss of $10\log_{10}(\pi/2)=1.96$\,dB (Chapter~\ref{ch:infotheory}); at rate $1/2$ the
hard-decision Shannon limit is about 1.8\,dB, against 0.19\,dB with soft decisions.

\mdcfig{ch14_hard_soft}{Hard versus soft decisions. Left: the Hamming (7,4) code decoded by
syndrome (hard) and by maximum-likelihood correlation against all 16 codewords (soft). Right: the
$K=7$ convolutional code with hard- and soft-decision Viterbi decoding. In both cases soft
decisions are worth between 1.5 and 2.5\,dB at error rates around $10^{-5}$.}

\begin{keyidea}[Never throw away information before the decoder]
A slicer is an irreversible data-processing step: it discards the reliability of each bit, and
no decoder can recover it. The rule that decoders should receive soft information, ideally
log-likelihood ratios, is one of the most valuable lessons in communications engineering, and it
applies with even more force to the iterative decoders of Chapter~\ref{ch:moderncodes}, which live
on soft information. In practice three or four bits of quantisation capture almost all of the
benefit (Section~\ref{sec:ch14:viterbi}).
\end{keyidea}

Why, then, did hard decisions survive so long? Because the classical algebraic decoders of the
following sections---for BCH and Reed--Solomon codes---work over finite fields and need hard
symbols. Soft-decision decoding of algebraic codes is possible (Chase algorithms, generalised
minimum-distance decoding, ordered-statistics decoding, and the Koetter--Vardy algebraic
soft-decision decoder for Reed--Solomon codes) but it is more complex and never displaced the
simple approach in mass-market systems. Convolutional codes, by contrast, were soft-decision from
the start, which is a large part of why they won the power-limited channels.

%======================================================================
\section{Finite fields made gentle}
\label{sec:ch14:gf}
%======================================================================

Hamming codes correct one error. To correct many, the codes of the rest of this chapter need
arithmetic in which we can add, subtract, multiply and \emph{divide} blocks of bits, solving
equations for the error positions just as we would with real numbers. The structure that
allows this is a \term{finite field}, or \term{Galois field}, named after \'Evariste Galois,
who discovered them in the 1830s shortly before dying in a duel at the age of twenty.

\subsection{Fields with a prime number of elements}

A \emph{field} is a set with addition and multiplication that behave like those of the rational
numbers: both are commutative and associative, multiplication distributes over addition, there
are identities 0 and 1, and every element has an additive inverse and every nonzero element a
multiplicative inverse. The integers modulo a prime $p$ form a field, GF($p$). In GF(2) the only
elements are 0 and 1, addition is XOR and multiplication is AND. In GF(7), $3\times5=15\equiv1$, so
$3^{-1}=5$. The integers modulo a composite number are \emph{not} a field: modulo 8,
$2\times4\equiv0$, so 2 has no inverse.

\subsection{GF($2^m$): polynomials modulo a primitive polynomial}

Digital hardware wants fields whose elements are $m$-bit words, and it turns out that a field with
$q$ elements exists if and only if $q$ is a prime power, and that it is unique up to relabelling.
GF($2^m$) is built exactly as the complex numbers are built from the reals. Take an
\term{irreducible polynomial} $p(x)$ of degree $m$ over GF(2)---one with no factors, just as
$x^2+1$ has no real factors---and invent a root $\alpha$ with $p(\alpha)=0$. The field elements are
the polynomials in $\alpha$ of degree less than $m$ with binary coefficients, which we store as
$m$-bit vectors:
\[
a_{m-1}\alpha^{m-1}+\dots+a_1\alpha+a_0\ \longleftrightarrow\ (a_{m-1}\dots a_1a_0).
\]
Addition is coefficient-wise modulo 2, that is, bitwise XOR. Multiplication is polynomial
multiplication followed by reduction using $p(\alpha)=0$. If, in addition, $p(x)$ is
\term{primitive}, then the powers $\alpha^0,\alpha^1,\dots,\alpha^{2^m-2}$ are all distinct and run
through every nonzero element, and $\alpha^{2^m-1}=1$. Every nonzero element is then a power of
$\alpha$, which gives two representations of each element: as a bit vector (convenient for
adding) and as a power, or \emph{logarithm} (convenient for multiplying).

Table~\ref{tab:ch14:gf8} builds GF(8) from $p(x)=x^3+x+1$. Starting from $\alpha^0=1$, each
row is the previous one multiplied by $\alpha$: a left shift of the bit vector, followed, whenever
an $\alpha^3$ term appears, by replacing it with $\alpha+1$ (since $\alpha^3=\alpha+1$ modulo $p$).
The same table for GF(16) with $p(x)=x^4+x+1$ is alongside.

\begin{table}[htbp]\centering\small
\caption{Log/antilog tables for GF(8) with $p(x)=x^3+x+1$ and GF(16) with $p(x)=x^4+x+1$. Each
nonzero element is a power of the primitive element $\alpha$.}
\label{tab:ch14:gf8}
\begin{tabular}{@{}c l c c@{}}
\toprule
\multicolumn{4}{c}{GF(8)}\\
power & polynomial & vector & int\\
\midrule
$0$ & $0$ & 000 & 0\\
$\alpha^0$ & $1$ & 001 & 1\\
$\alpha^1$ & $\alpha$ & 010 & 2\\
$\alpha^2$ & $\alpha^2$ & 100 & 4\\
$\alpha^3$ & $\alpha+1$ & 011 & 3\\
$\alpha^4$ & $\alpha^2+\alpha$ & 110 & 6\\
$\alpha^5$ & $\alpha^2+\alpha+1$ & 111 & 7\\
$\alpha^6$ & $\alpha^2+1$ & 101 & 5\\
\bottomrule
\end{tabular}\qquad
\begin{tabular}{@{}c c c c c c@{}}
\toprule
\multicolumn{6}{c}{GF(16)}\\
power & vector & int & power & vector & int\\
\midrule
$\alpha^0$ & 0001 & 1 & $\alpha^8$ & 0101 & 5\\
$\alpha^1$ & 0010 & 2 & $\alpha^9$ & 1010 & 10\\
$\alpha^2$ & 0100 & 4 & $\alpha^{10}$ & 0111 & 7\\
$\alpha^3$ & 1000 & 8 & $\alpha^{11}$ & 1110 & 14\\
$\alpha^4$ & 0011 & 3 & $\alpha^{12}$ & 1111 & 15\\
$\alpha^5$ & 0110 & 6 & $\alpha^{13}$ & 1101 & 13\\
$\alpha^6$ & 1100 & 12 & $\alpha^{14}$ & 1001 & 9\\
$\alpha^7$ & 1011 & 11 & $\alpha^{15}$ & $=\alpha^0$ & \\
\bottomrule
\end{tabular}
\end{table}

\begin{worked}[Arithmetic in GF(8)]
Add $\alpha^4$ and $\alpha^5$: in vector form $110\oplus111=001$, so $\alpha^4+\alpha^5=1$.
Multiply them: add the logarithms modulo 7, $\alpha^4\alpha^5=\alpha^9=\alpha^2=100$. Check by
polynomials: $(\alpha^2+\alpha)(\alpha^2+\alpha+1)=\alpha^4+\alpha$ (the cross terms cancel in
pairs), and $\alpha^4=\alpha^2+\alpha$, so the product is $\alpha^2$. Divide: $\alpha^2/\alpha^5=
\alpha^{-3}=\alpha^4$. Invert: $(\alpha^3)^{-1}=\alpha^{4}$, and indeed $011\times110$ reduces to 001.
Note that $x+x=0$ for every $x$: in a field of characteristic 2, addition and subtraction are the
same operation, which removes every sign from the decoding algorithms that follow.
\end{worked}

\subsection{Conjugates and minimal polynomials}

Binary codes need one more idea. If $\beta$ is a root of a polynomial with \emph{binary}
coefficients, so is $\beta^2$, because squaring is linear in characteristic 2:
$(a+b)^2=a^2+b^2$. The elements $\beta,\beta^2,\beta^4,\dots$ are the \term{conjugates} of $\beta$,
and the lowest-degree binary polynomial having $\beta$ as a root, the \term{minimal polynomial}
$m_\beta(x)$, is the product of $(x-\gamma)$ over all conjugates $\gamma$. In GF(16) the exponents
split into the cyclotomic cosets $\{1,2,4,8\}$, $\{3,6,12,9\}$, $\{5,10\}$ and $\{7,14,13,11\}$,
with minimal polynomials
\[
m_1=x^4+x+1,\quad m_3=x^4+x^3+x^2+x+1,\quad m_5=x^2+x+1,\quad m_7=x^4+x^3+1 .
\]
These are the building blocks of BCH generator polynomials (Section~\ref{sec:ch14:bch}).

\subsection{Implementation}

Hardware adds field elements with XOR gates and multiplies them with small arrays of AND and XOR
gates (a general GF(256) multiplier needs 64 AND gates and roughly as many XOR gates);
multiplication by a \emph{constant}, which is what encoders need, is cheaper still. Software uses
log and antilog tables: $ab=\exp[(\log a+\log b)\bmod(2^m-1)]$, with a special case for zero, or a
full $256\times256$ product table for GF(256). Modern processors accelerate both: carry-less
multiply instructions (PCLMULQDQ on x86, PMULL on Arm) compute the unreduced polynomial product,
and the x86 GFNI extension multiplies bytes in GF(256) directly.

\begin{pitfall}[Irreducible is not primitive, and fields are not interchangeable]
All GF(256)s are isomorphic, but their bit representations differ with the choice of $p(x)$, and
codes built over one are incompatible with another. DVB, ATSC, QR codes, CD audio and RAID-6 use
$x^8+x^4+x^3+x^2+1$ (\texttt{0x11D}); CCSDS uses $x^8+x^7+x^2+x+1$ (\texttt{0x187}) together with a
``dual basis'' representation of the bytes; AES uses $x^8+x^4+x^3+x+1$ (\texttt{0x11B}), which is
irreducible but \emph{not} primitive: $x$ has order 51 rather than 255, and $x+1$ must be used as the
generator of log tables. A decoder built with the wrong polynomial, the wrong generator or the
wrong bit order will fail on every codeword with a nonzero syndrome, which is a maddening bug
because it passes every error-free test vector.
\end{pitfall}

Table~\ref{tab:ch14:prim} lists the primitive polynomials used in practice. The list in
\lab{commlib/gf.py} matches it.

\begin{table}[htbp]\centering\small
\caption{Commonly used primitive polynomials for GF($2^m$). Hex form includes the $x^m$ term.}
\label{tab:ch14:prim}
\begin{tabularx}{\linewidth}{@{}c l l X@{}}
\toprule
$m$ & polynomial & hex & typical use\\
\midrule
3 & $x^3+x+1$ & \texttt{0xB} & textbook examples, RS(7,3)\\
4 & $x^4+x+1$ & \texttt{0x13} & BCH(15,$k$), CRC-4 building block\\
6 & $x^6+x+1$ & \texttt{0x43} & BCH(63,$k$)\\
8 & $x^8+x^4+x^3+x^2+1$ & \texttt{0x11D} & RS in CD, DVB, ATSC, QR, RAID-6\\
8 & $x^8+x^7+x^2+x+1$ & \texttt{0x187} & CCSDS RS(255,223)\\
10 & $x^{10}+x^3+1$ & \texttt{0x409} & RS(528,514) and RS(544,514) in 100G--800G Ethernet\\
14 & $x^{14}+x^5+x^3+x+1$ & \texttt{0x402B} & DVB-S2 BCH, short frames\\
16 & $x^{16}+x^5+x^3+x^2+1$ & \texttt{0x1002D} & DVB-S2 BCH, normal frames\\
\bottomrule
\end{tabularx}
\end{table}

%======================================================================
\section{Cyclic codes}
\label{sec:ch14:cyclic}
%======================================================================

A linear code is \term{cyclic} if every cyclic shift of a codeword is also a codeword. The
definition looks like a curiosity, but it gives codes an algebraic structure that makes encoding
a shift register and decoding a matter of polynomial arithmetic. Hamming, Golay, BCH and
Reed--Solomon codes are all cyclic (or shortened cyclic), and so is every CRC.

\subsection{The polynomial view}

Write a word $(c_{n-1},\dots,c_1,c_0)$ as the polynomial $c(x)=c_{n-1}x^{n-1}+\dots+c_1x+c_0$.
Multiplying by $x$ shifts every coefficient up one place and pushes $c_{n-1}$ out to $x^n$;
reducing modulo $x^n-1$ (that is, replacing $x^n$ by 1) wraps it round to the constant term. So a
cyclic shift is multiplication by $x$ modulo $x^n-1$, and a cyclic code is a set of polynomials
closed under addition and under multiplication by $x$, hence by any polynomial. The consequences
follow quickly:
\begin{itemize}
\item Every cyclic code has a unique monic \term{generator polynomial} $g(x)$ of lowest degree,
of degree $n-k$, and every codeword is a multiple of it: $c(x)=u(x)g(x)$ with $\deg u<k$.
\item $g(x)$ divides $x^n-1$. Conversely, every divisor of $x^n-1$ generates a cyclic code. The
codes of length $n$ are therefore catalogued by the factorisation of $x^n-1$.
\item The \term{parity-check polynomial} $h(x)=(x^n-1)/g(x)$ satisfies $c(x)h(x)\equiv0\bmod
(x^n-1)$ for every codeword.
\end{itemize}
Over GF(2), $x^7-1=(x+1)(x^3+x+1)(x^3+x^2+1)$. The generator $g(x)=x^3+x+1$ gives a (7,4) cyclic
code, which is a Hamming code with its columns reordered; $g(x)=(x+1)(x^3+x+1)$ gives the (7,3)
code of its even-weight words; $g(x)=x+1$ gives the (7,6) single-parity-check code.

\subsection{Systematic encoding by division}

The product $u(x)g(x)$ is a valid but non-systematic codeword. The systematic encoder instead
shifts the message into the high-order positions and fills the low-order positions with the
remainder of a division:
\begin{equation}
c(x)=x^{n-k}u(x)+\rho(x),\qquad \rho(x)=x^{n-k}u(x)\bmod g(x).
\label{eq:ch14:sysenc}
\end{equation}
Since $x^{n-k}u(x)=q(x)g(x)+\rho(x)$ and addition equals subtraction, $c(x)=q(x)g(x)$ is a
multiple of $g(x)$, hence a codeword, and its top $k$ coefficients are the message. The division
is performed by a \term{linear-feedback shift register} (LFSR) with $n-k$ stages and taps set by
the coefficients of $g(x)$ (Figure~\ref{fig:ch14_lfsr}). The message bits are clocked in, most
significant first; after the last one, the register holds $\rho(x)$, which is shifted out behind
the message. The circuit needs $n-k$ flip-flops and at most $n-k$ XOR gates, whatever the block
length.

\begin{figure}[htbp]\centering
\begin{tikzpicture}[>=Stealth,
  reg/.style={draw=navy,thick,fill=navy!6,minimum width=0.9cm,minimum height=0.75cm,font=\sffamily\small},
  xr/.style={draw=accent,thick,circle,inner sep=0pt,minimum size=0.45cm,font=\small},
  lab/.style={font=\sffamily\scriptsize,text=navy}]
\node[xr] (x0) at (0,0) {$+$};
\node[reg] (r0) at (1.3,0) {$r_0$};
\node[xr] (x1) at (2.6,0) {$+$};
\node[reg] (r1) at (3.9,0) {$r_1$};
\node[reg] (r2) at (5.4,0) {$r_2$};
\node[reg] (r3) at (6.9,0) {$r_3$};
\node[xr] (xin) at (8.3,0) {$+$};
\draw[->,thick,navy] (x0) -- (r0); \draw[->,thick,navy] (r0) -- (x1); \draw[->,thick,navy] (x1) -- (r1);
\draw[->,thick,navy] (r1) -- (r2); \draw[->,thick,navy] (r2) -- (r3); \draw[->,thick,navy] (r3) -- (xin);
\draw[->,thick,navy] (9.6,0) node[right,lab,align=left]{message\\$u_{k-1},\dots,u_0$} -- (xin);
\draw[thick,accent] (xin) -- ++(0,-1.1) coordinate (fb);
\draw[->,thick,accent] (fb) -| (x0);
\draw[->,thick,accent] (x1|-fb) -- (x1);
\node[lab,above=0.05cm] at ($(fb)!0.3!(x1|-fb)$) {feedback};
\node[lab] at (0,0.6) {$g_0=1$}; \node[lab] at (2.6,0.6) {$g_1=1$};
\node[lab] at (4.65,0.6) {$g_2=0$}; \node[lab] at (6.15,0.6) {$g_3=0$};
\node[lab,align=center] at (8.3,0.75) {$g_4=1$};
\node[lab,align=left,anchor=north west] at (-0.3,-1.45) {After the last message bit, $r_3r_2r_1r_0$ holds the remainder
$\rho(x)=x^4u(x)\bmod(x^4+x+1)$,\\which is shifted out after the message. Clocking a received word through the same
circuit\\(without the premultiplication) leaves its syndrome.};
\end{tikzpicture}
\caption{A systematic encoder for the cyclic code, or CRC, with generator $g(x)=x^4+x+1$. Each
clock computes the feedback bit (the incoming message bit plus the register's top bit) and adds it
into the stages where $g(x)$ has a nonzero coefficient; the XOR at each tap performs one step of
polynomial long division.}
\label{fig:ch14_lfsr}
\end{figure}

\subsection{Syndromes and bursts}

The receiver computes $s(x)=r(x)\bmod g(x)$ with the same circuit. The syndrome is zero if and
only if $r(x)$ is a codeword, and it equals $e(x)\bmod g(x)$, so it depends only on the error
pattern. Cyclic structure gives a strong guarantee against \term{burst errors}, error patterns
confined to $b$ consecutive positions. Such a pattern can be written $e(x)=x^ib(x)$ with $\deg
b(x)<b$ and $b(0)=1$. If $b\le n-k$, then $b(x)$ has lower degree than $g(x)$ and so is not a
multiple of it, and $x^i$ shares no factor with $g(x)$ (whose constant term is 1), so $e(x)$ is
not a codeword and the error is detected:
\begin{quote}
\emph{A cyclic or shortened cyclic code with $n-k$ check bits detects every burst of length
$n-k$ or less.} Of the bursts of length $n-k+1$ a fraction $2^{-(n-k-1)}$ go undetected, and of
longer bursts a fraction $2^{-(n-k)}$.
\end{quote}
The same structure yields simple \emph{correcting} decoders for single errors and short bursts:
the Meggitt decoder and error trapping shift the received word cyclically through the syndrome
register and watch for a syndrome pattern that identifies a correctable error at a fixed
position. Fire codes, designed for burst correction on disks and tapes in the 1960s, used exactly
this approach.

%======================================================================
\section{The cyclic redundancy check}
\label{sec:ch14:crc}
%======================================================================

A \term{cyclic redundancy check} (CRC) is a shortened cyclic code used only for error
\emph{detection}. The transmitter appends $r$ check bits computed by~\eqref{eq:ch14:sysenc} to a
message of arbitrary length; the receiver recomputes them and discards the frame on a mismatch.
W.~Wesley Peterson and D.~T.~Brown described the technique in 1961, and it now protects nearly
every frame in every network: Ethernet, Wi-Fi, USB, Bluetooth, PCI Express, SATA, the transport
and code blocks of LTE and 5G, the ZIP and PNG file formats, and the file systems on your disk.

\begin{worked}[A CRC by hand]
Compute the 4-bit CRC of the message $1101011011$ with $g(x)=x^4+x+1$ (binary $10011$).

Append four zeros (multiply by $x^4$) and divide by $10011$ with XOR in place of subtraction,
writing the divisor under each leading one:
\begin{center}\small
\begin{tabular}{@{}l@{\hspace{2em}}l@{}}
\texttt{11010110110000} & dividend $x^4u(x)$\\
\texttt{10011}\phantom{\texttt{000000000}} & $\to$ \texttt{01001110110000}\\
\texttt{\phantom{1}10011}\phantom{\texttt{00000000}} & $\to$ \texttt{00000010110000}\\
\texttt{\phantom{110101}10011}\phantom{\texttt{000}} & $\to$ \texttt{00000000101000}\\
\texttt{\phantom{11010110}10011}\phantom{\texttt{0}} & $\to$ \texttt{00000000001110}\\
\end{tabular}
\end{center}
The remainder is $1110$, so the transmitted frame is $1101011011\,1110$. Dividing the whole frame
by $10011$ leaves remainder zero. If the channel flips the third bit, the receiver's remainder is
$x^{11}\bmod g(x)$: since $x^4=x+1$, $x^8=x^2+1$ and $x^{11}=x^3\cdot x^8=x^5+x^3=x^3+x^2+x$, the
remainder is $1110$, which is nonzero, and the error is detected. The LFSR of
Figure~\ref{fig:ch14_lfsr} computes exactly the same remainder, one bit per clock.
\end{worked}

\subsection{What a CRC detects}

The detection properties of a CRC follow from the factors of its generator.
\begin{itemize}
\item \textbf{All single errors}, since $g(x)$ has at least two terms.
\item \textbf{All odd numbers of errors} if $(x+1)$ divides $g(x)$, since then every codeword has
$c(1)=0$, an even number of ones. Most standard CRCs include this factor.
\item \textbf{All double errors} within the order of $g(x)$: an error $x^i+x^j=x^i(x^{j-i}+1)$ is
missed only if $g(x)$ divides $x^{j-i}+1$. If $g(x)=(x+1)p(x)$ with $p(x)$ primitive of degree
$r-1$, this cannot happen for frames shorter than $2^{r-1}-1$ bits.
\item \textbf{All bursts of length $r$ or less}, and all but a fraction of about $2^{-r}$ of longer
ones, as shown above.
\end{itemize}
For completely random corruption, the probability that a frame passes the check is essentially
$2^{-r}$: one in 4.3 billion for a 32-bit CRC. For a binary symmetric channel with a low error
rate, the probability of undetected error is dominated by the lowest-weight codewords of the
shortened code, $P_{\text{ud}}\approx A_{d}\,p^{d}$, and the \emph{Hamming distance} $d$ of the CRC
at the frame length of interest becomes the key figure of merit. Philip Koopman's exhaustive
searches of the early 2000s showed that many long-standardised polynomials are far from the best
available. The IEEE 802.3 CRC-32, for example, has Hamming distance 4 over any frame length up to
about 91\,000 bits (covering every Ethernet frame), 5 up to about 3000 bits and 6 only up to about
270 bits, while the Castagnoli polynomial (CRC-32C) keeps distance 6 out to roughly 5000 bits.
CRC-32C was accordingly adopted for iSCSI, SCTP and several file systems.

Figure~\ref{fig:ch14_crc_pud} computes $P_{\text{ud}}$ exactly for some short CRCs, using the
MacWilliams identity: the dual of an $r$-bit CRC code has only $2^r$ codewords, whose weights can be
enumerated directly, and
\begin{equation}
P_{\text{ud}}(p)=2^{-r}\sum_{j=0}^{n}B_j(1-2p)^j-(1-p)^n,
\label{eq:ch14:macwilliams}
\end{equation}
where $B_j$ is the dual weight distribution and $n$ the frame length. Good polynomials have
$P_{\text{ud}}$ rising smoothly to $2^{-r}$ as $p\to\tfrac12$; the poorly chosen $x^8+1$ misses
every pair of errors eight bits apart, and at high error rates its undetected-error probability
even exceeds $2^{-r}$.

\mdcfig{ch14_crc_pud}{Probability of undetected error of CRCs on a binary symmetric channel,
computed exactly from~\eqref{eq:ch14:macwilliams}. Left: two good 8-bit polynomials (Hamming
distance 4 at this length, slope 4 on the log--log plot) and a poor one (distance 2). Right: the
16-bit CCITT polynomial at three message lengths, and the CRC-16/ARC polynomial. Longer messages
contain more low-weight codewords, so $P_{\text{ud}}$ rises with length at low $p$.}

\subsection{The standard CRCs}

Table~\ref{tab:ch14:crc} collects the polynomials an engineer meets most often. The hexadecimal
form is the usual ``normal'' representation, with the leading $x^r$ term omitted.

\begin{table}[htbp]\centering\small
\caption{Widely used CRC polynomials.}
\label{tab:ch14:crc}
\begin{tabularx}{\linewidth}{@{}l l X@{}}
\toprule
Name & Poly (hex) & Where\\
\midrule
CRC-6 (NR) & \texttt{0x21} & $D^6+D^5+1$; 5G NR uplink control information on polar codes (TS~38.212)\\
CRC-8/ATM & \texttt{0x07} & ATM cell header error control (HEC); also corrects single errors\\
CRC-8 (LTE) & \texttt{0x9B} & LTE CQI/control reports (TS~36.212)\\
CRC-11 (NR) & \texttt{0x621} & 5G NR uplink control information (TS~38.212)\\
CRC-16-CCITT & \texttt{0x1021} & HDLC, X.25, Bluetooth, SD cards; NR transport blocks up to 3824 bits\\
CRC-16/IBM (ARC) & \texttt{0x8005} & USB data packets (with inversion), Modbus, legacy protocols\\
CRC-24A (LTE/NR) & \texttt{0x864CFB} & Transport-block CRC (TS~36.212, TS~38.212)\\
CRC-24B (LTE/NR) & \texttt{0x800063} & Code-block CRC for turbo and LDPC segments\\
CRC-24C (NR) & \texttt{0xB2B117} & NR downlink control and broadcast channels (polar, distributed CRC)\\
CRC-32 & \texttt{0x04C11DB7} & Ethernet, Wi-Fi, SATA, ZIP, PNG, MPEG-2 TS sections\\
CRC-32C & \texttt{0x1EDC6F41} & iSCSI, SCTP, ext4 and Btrfs metadata; SSE4.2 \texttt{crc32} instruction\\
CRC-64 (ECMA-182) & \texttt{0x42F0E1EBA9EA3693} & tape and storage formats, XZ archives\\
\bottomrule
\end{tabularx}
\end{table}

\begin{pitfall}[A CRC is a polynomial plus five conventions]
Two implementations of ``CRC-32'' with the same polynomial will disagree unless they also agree
on: the \emph{initial value} of the register (all ones for CRC-32, so that leading zero bytes
change the result); whether input bytes are \emph{reflected} (processed least significant bit
first, as serial hardware sends them); whether the final remainder is reflected; the value
\emph{XORed onto the output} (all ones for CRC-32); and the order in which the check bytes are
appended. Ross Williams's 1993 ``painless guide'' introduced the parameterisation (width, poly,
init, refin, refout, xorout) used by every modern CRC catalogue, together with a \emph{check}
value: the CRC of the ASCII string \texttt{123456789}, which is \texttt{0xCBF43926} for CRC-32. Test
against the check value first, and most CRC bugs disappear. \lab{commlib/gf.py} includes this
parameterised CRC with several catalogue entries.
\end{pitfall}

\begin{inpractice}[Fast CRCs]
The bit-serial LFSR processes one bit per clock, which was fine for 1970s modems but not for a
100\,Gb/s link. Because the CRC is linear, the register's next state after a whole byte is a
linear function of its current state and the byte, which can be precomputed. Sarwate's 1988
table-driven algorithm uses a 256-entry table to process a byte per lookup; ``slicing-by-8''
uses eight tables to process eight bytes per iteration. Hardware designers unroll the same
linear map into XOR trees that process 64 to 1024 bits per clock. Modern x86 processors compute
CRC-32C with a dedicated instruction and fold long buffers with the carry-less multiply
instruction at many gigabytes per second per core.
\end{inpractice}

\begin{inpractice}[CRCs as the referee of modern decoders]
In LTE every transport block carries a 24-bit CRC (in 5G NR, 16 bits for small blocks), and every
code block of a segmented transport block carries its own. The CRC decides the HARQ acknowledgement
(Section~\ref{sec:ch14:arq}) and lets an LDPC or turbo decoder stop iterating as soon as the
codeword checks. In the polar codes of the 5G control channels, the CRC is part of the code: the
successive-cancellation list decoder produces several candidate messages and the CRC picks the
one that is consistent (Chapter~\ref{ch:moderncodes}). The same CRC is scrambled with the
identity of the intended terminal, so a phone detects both errors and ``not addressed to me'' with
one check.
\end{inpractice}

%======================================================================
\section{BCH codes}
\label{sec:ch14:bch}
%======================================================================

Hamming codes correct one error. In 1959 Alexis Hocquenghem, and independently in 1960 Raj Chandra
Bose and Dijen Ray-Chaudhuri, found how to design cyclic codes that correct any prescribed number
of errors $t$. The \term{BCH codes} named after them remain among the best known codes of short
and moderate length, and they can be decoded with algebraic algorithms whose complexity grows only
polynomially with $t$.

\subsection{Designing for distance: the BCH bound}

The key is to specify a cyclic code by the \emph{roots} of its generator polynomial in an
extension field. Let $\alpha$ be a primitive element of GF($2^m$) and $n=2^m-1$. If $g(x)$ has the
$2t$ consecutive powers $\alpha,\alpha^2,\dots,\alpha^{2t}$ among its roots, every codeword
satisfies $c(\alpha^j)=0$ for $j=1,\dots,2t$, and the code has minimum distance at least $2t+1$.
The proof is short and worth knowing. Suppose a codeword had weight $w\le2t$, with nonzero bits at
positions $i_1,\dots,i_w$. The conditions $c(\alpha^j)=\sum_l\alpha^{ji_l}=0$ for $j=1,\dots,w$
form a homogeneous linear system in the $w$ unknown bits whose matrix is a Vandermonde matrix in the
distinct elements $\alpha^{i_l}$, which is nonsingular. So the only solution is zero, a
contradiction. This is the \term{BCH bound}: $\delta-1$ consecutive roots guarantee distance at
least $\delta$, the \emph{designed distance}.

Since $g(x)$ must have binary coefficients, whenever $\alpha^j$ is a root so are all its
conjugates $\alpha^{2j},\alpha^{4j},\dots$. The generator is therefore the least common multiple
of minimal polynomials,
\begin{equation}
g(x)=\operatorname{lcm}\{m_1(x),m_3(x),\dots,m_{2t-1}(x)\},
\label{eq:ch14:bchgen}
\end{equation}
where only the odd indices are needed because $m_{2j}=m_j$. Each minimal polynomial has degree at
most $m$, so $n-k\le mt$: a $t$-error-correcting BCH code of length $2^m-1$ needs at most $mt$ check
bits. Compare the Hamming bound, which requires about $t\log_2 n\approx mt$ check bits: for small
$t$ the BCH codes are close to optimal, which is what the right panel of Figure~\ref{fig:ch14_bounds}
shows.

For $m=4$, using the minimal polynomials of Section~\ref{sec:ch14:gf}: $t=1$ gives
$g=m_1=x^4+x+1$, the (15,11) Hamming code; $t=2$ gives $g=m_1m_3=x^8+x^7+x^6+x^4+1$, a (15,7)
double-error-correcting code; $t=3$ gives $g=m_1m_3m_5=x^{10}+x^8+x^5+x^4+x^2+x+1$, a (15,5)
code correcting three errors.

\subsection{Decoding: syndromes, error locators and the key idea}

Suppose the received word has $\nu\le t$ errors at positions $i_1,\dots,i_\nu$. Call
$X_l=\alpha^{i_l}$ the \term{error locators}. The decoder computes $2t$ syndromes by evaluating the
received polynomial at the roots of $g(x)$:
\begin{equation}
S_j=r(\alpha^j)=e(\alpha^j)=\sum_{l=1}^{\nu}X_l^{\,j},\qquad j=1,\dots,2t .
\label{eq:ch14:bchsyn}
\end{equation}
These are $2t$ nonlinear equations (power sums) in the $\nu$ unknown locators. The trick that
makes them solvable is to introduce the \term{error-locator polynomial}
\begin{equation}
\Lambda(x)=\prod_{l=1}^{\nu}(1-X_lx)=1+\Lambda_1x+\dots+\Lambda_\nu x^\nu,
\label{eq:ch14:locator}
\end{equation}
whose roots are the \emph{inverses} of the locators. Its coefficients are related to the syndromes
by Newton's identities, which in characteristic 2 take the form of a linear recurrence:
\begin{equation}
S_j+\Lambda_1S_{j-1}+\dots+\Lambda_\nu S_{j-\nu}=0,\qquad j=\nu+1,\dots,2t .
\label{eq:ch14:newton}
\end{equation}
The syndromes are the output of a linear-feedback shift register whose connection polynomial is
$\Lambda(x)$. Decoding becomes a sequence of three steps:
\begin{enumerate}
\item compute the syndromes~\eqref{eq:ch14:bchsyn};
\item find $\Lambda(x)$, the shortest LFSR that generates them;
\item find the roots of $\Lambda(x)$ by trying every field element (the \term{Chien search}),
and flip the bits at the corresponding positions.
\end{enumerate}
Peterson (1960) solved step 2 by inverting the $\nu\times\nu$ linear system~\eqref{eq:ch14:newton},
trying $\nu=t,t-1,\dots$ until the matrix is nonsingular, which costs $O(t^3)$ operations. Berlekamp's
algorithm of 1968, which Massey reinterpreted in 1969 as the synthesis of the shortest LFSR generating
a given sequence, does it in $O(t^2)$ and is the standard today. For $t=2$ the answer is explicit:
$\Lambda_1=S_1$ and $\Lambda_2=(S_3+S_1^3)/S_1$.

\begin{worked}[Decoding the (15,7) BCH code]
The all-zero codeword is sent and errors occur at positions 3 and 10, so $e(x)=x^{10}+x^3$. Using
the GF(16) table: $S_1=\alpha^{10}+\alpha^3=0111\oplus1000=1111=\alpha^{12}$, and
$S_3=\alpha^{30}+\alpha^9=\alpha^0+\alpha^9=0001\oplus1010=1011=\alpha^7$.

Then $\Lambda_1=S_1=\alpha^{12}$ and $S_1^3=\alpha^{36}=\alpha^{6}$, so
$\Lambda_2=(\alpha^7+\alpha^6)/\alpha^{12}=(1011\oplus1100)/\alpha^{12}=\alpha^{10}/\alpha^{12}
=\alpha^{13}$. The locator is $\Lambda(x)=1+\alpha^{12}x+\alpha^{13}x^2$. Check it with the true
locators: $X_1+X_2=\alpha^3+\alpha^{10}=\alpha^{12}$ and $X_1X_2=\alpha^{13}$, as required.

The Chien search evaluates $\Lambda(\alpha^{-i})$ for $i=0,\dots,14$. At $i=3$,
$\Lambda(\alpha^{12})=1+\alpha^{24}+\alpha^{37}=1+\alpha^9+\alpha^7=0001\oplus1010\oplus1011=0$; at
$i=10$, $\Lambda(\alpha^5)=1+\alpha^{17}+\alpha^{23}=1+\alpha^2+\alpha^8=0001\oplus0100\oplus0101=0$.
No other position gives zero. The decoder flips bits 3 and 10 and the codeword is restored.
\end{worked}

\subsection{BCH codes in practice}

BCH codes are the code of choice when a moderate number of \emph{random} bit errors must be
corrected with a simple, deterministic, hard-decision decoder, and the list of applications is
long. The POCSAG paging protocol protects each 32-bit word with a (31,21) BCH code and a parity
bit. Many NAND flash controllers of the 2000s and early 2010s used BCH codes over GF($2^{13}$) or
GF($2^{14}$) correcting from 4 to about 40 errors per 512-byte or 1-kilobyte sector, until the
rising raw error rates of multi-level cells forced a move to soft-decision LDPC codes. And in the
DVB-S2 satellite standard (ETSI EN~302~307), every LDPC codeword is preceded by an outer BCH code
correcting 8, 10 or 12 errors over GF($2^{16}$), whose job is to clean up the rare residual errors
of the LDPC decoder (the ``error floor''), a small-overhead insurance policy that is typical of
how the classical codes survive in modern systems.

\begin{inpractice}[DVB-S2: BCH as an outer code]
A DVB-S2 normal frame carries 64\,800 coded bits. The LDPC inner code has rates from 1/4 to 9/10;
the BCH outer code adds 128, 160 or 192 parity bits ($t=8$, 10 or 12) to the data. At rate 1/2,
for example, the BCH code maps 32\,208 data bits to 32\,400, an overhead of about 0.6\%. The LDPC
decoder does almost all the work; the BCH decoder removes the few isolated errors the LDPC decoder
occasionally leaves behind, so that the quasi-error-free target (a packet error rate of about
$10^{-7}$) is met without the LDPC code having to reach it alone. DVB-T2, DVB-C2 and ATSC~3.0
adopted the same BCH-plus-LDPC structure.
\end{inpractice}

%======================================================================
\section{Reed--Solomon codes}
\label{sec:ch14:rs}
%======================================================================

If one code from this chapter had to be singled out as the most successful, it would be the
Reed--Solomon code. It is in every CD, DVD and Blu-ray disc, every QR code, most digital television
standards, deep-space telemetry, RAID-6 disk arrays, cloud storage, hard-disk and tape sectors, DSL
modems and the forward error correction of 100 to 800\,Gb/s Ethernet. The reason is that it is
simultaneously optimal in distance (MDS), naturally matched to bytes, superb against bursts, and
decodable with algorithms that fit in a small amount of silicon.

\subsection{Two definitions}

Irving Reed and Gustave Solomon, at MIT Lincoln Laboratory, defined their codes in 1960 by
\emph{evaluation}. Take a field GF($q$) and $n\le q$ distinct field elements
$\beta_1,\dots,\beta_n$. Treat the $k$ message symbols as the coefficients of a polynomial $m(x)$ of
degree less than $k$, and transmit its values:
\begin{equation}
\vect{c}=\big(m(\beta_1),m(\beta_2),\dots,m(\beta_n)\big).
\label{eq:ch14:rseval}
\end{equation}
Two different polynomials of degree less than $k$ can agree in at most $k-1$ points (their
difference has at most $k-1$ roots), so two different codewords differ in at least $n-k+1$
positions. The code meets the Singleton bound: Reed--Solomon codes are MDS. The same argument
shows the MDS property directly: any $k$ received values determine $m(x)$ by interpolation, so any
$k$ of the $n$ symbols suffice to recover the message. That is the essence of erasure coding.

The now-standard definition, with $n=q-1$ and $\beta_i=\alpha^{i}$, makes the code cyclic with
generator
\begin{equation}
g(x)=\prod_{j=0}^{n-k-1}\big(x-\alpha^{b+j}\big),
\label{eq:ch14:rsgen}
\end{equation}
for some integer $b$ (the ``first consecutive root'', 0 or 1 in most standards; 112 in CCSDS).
Here the coefficients of $g(x)$ are field elements, not bits: a Reed--Solomon code is a BCH-type
code over GF($2^m$) itself, and each minimal polynomial is just the linear factor $x-\alpha^j$. With
$n-k=2t$ consecutive roots the BCH bound gives $d\ge2t+1$, and Singleton gives $d\le2t+1$, so
$d=n-k+1$ exactly. Encoding is~\eqref{eq:ch14:sysenc} with field arithmetic: an LFSR with $2t$
byte-wide stages and constant multipliers.

\subsection{Why symbols beat bits against bursts}

A Reed--Solomon code over GF(256) corrects $t$ \emph{symbol} errors, and a symbol is in error
whether one of its bits or all eight are wrong. A burst of $B$ bit errors touches at most
$\lceil(B-1)/8\rceil+1$ bytes, so the RS(255,223) code with $t=16$ corrects any single burst of up
to $8\times15+1=121$ bits, and, if the bursts are flagged as erasures, twice that. A binary BCH code
of about the same length and rate corrects only about 23 bit errors, so a single 24-bit burst defeats it. On channels whose errors
come in clusters---a scratch on a disc, a fade on a radio link, the error bursts that leave a
Viterbi decoder---symbol-level coding is the natural fit.

The right panel of Figure~\ref{fig:ch14_rs_performance} shows the other face of RS codes. Output
error rate falls extremely steeply once the input symbol error rate drops below about $t/n$: a
long, powerful code such as RS(255,223) goes from useless to nearly perfect over a narrow range
of input error rates, the \term{cliff effect} familiar to anyone who has watched digital
television break up. The left panel shows the corresponding bit error rates with BPSK and hard
decisions.

\mdcfig{ch14_rs_performance}{Reed--Solomon performance. Left: bit error rate of several
byte-oriented codes with BPSK and hard decisions, assuming independent bit errors. RS(204,188), the
shortened code of DVB, performs almost exactly like its parent RS(255,239). Right: decoded symbol
error rate against channel symbol error rate for $n=255$ and $t$ from 1 to 16: the larger $t$,
the sharper the threshold.}

\subsection{Decoding Reed--Solomon codes}

The decoder follows the BCH recipe with one extra step: because symbols are not bits, the decoder
must find the error \emph{values} as well as their locations. With error locators $X_l$ and
values $Y_l$, the syndromes are
\begin{equation}
S_j=r(\alpha^{b+j})=\sum_{l=1}^{\nu}Y_l\,X_l^{\,b+j},\qquad j=0,\dots,2t-1.
\end{equation}
Define the syndrome polynomial $S(x)=\sum_jS_jx^j$. The error locator $\Lambda(x)$
of~\eqref{eq:ch14:locator} and an \term{error-evaluator polynomial} $\Omega(x)$ of degree less than
$\nu$ satisfy the \term{key equation}
\begin{equation}
S(x)\,\Lambda(x)\equiv\Omega(x)\pmod{x^{2t}}.
\label{eq:ch14:key}
\end{equation}
The complete decoder has four steps:
\begin{enumerate}
\item \textbf{Syndromes.} Evaluate $r(x)$ at the $2t$ roots of $g(x)$ by Horner's rule: $2tn$
multiply-accumulates. If all are zero, accept the word.
\item \textbf{Key equation.} Find $\Lambda(x)$ with the Berlekamp--Massey algorithm (or Sugiyama's
1975 variant of Euclid's algorithm, popular in hardware), then $\Omega(x)=S(x)\Lambda(x)\bmod
x^{2t}$.
\item \textbf{Chien search.} Evaluate $\Lambda(\alpha^{-i})$ for every position $i$; each zero marks
an error at position $i$. If the number of roots differs from $\deg\Lambda$, more than $t$ errors
occurred: declare a decoding failure.
\item \textbf{Forney's formula.} Compute each error value as
\begin{equation}
Y_l=X_l^{\,1-b}\,\frac{\Omega(X_l^{-1})}{\Lambda'(X_l^{-1})},
\label{eq:ch14:forney}
\end{equation}
where $\Lambda'$ is the formal derivative (in characteristic 2, just the odd-power terms of
$\Lambda$ shifted down), and add it to the received symbol.
\end{enumerate}
The Berlekamp--Massey algorithm itself is a dozen lines. It processes the syndromes one at a time,
maintaining the current shortest LFSR $\Lambda(x)$ and its length $L$. At step $r$ it computes the
\emph{discrepancy} $\Delta=S_r+\sum_{i=1}^{L}\Lambda_iS_{r-i}$, the difference between the next
syndrome and the LFSR's prediction. If $\Delta=0$ the LFSR is still correct. Otherwise it is
corrected by adding a scaled, shifted copy of the LFSR saved at the last length change,
$\Lambda(x)\leftarrow\Lambda(x)-(\Delta/\Delta_{\text{old}})\,x^{s}B(x)$, and if $2L\le r$ the length
increases to $r+1-L$. After $2t$ steps $\Lambda(x)$ is the error locator. The implementation in
\lab{commlib/gf.py} follows exactly this outline and is worth reading alongside the text.

\begin{worked}[A Reed--Solomon code over GF(8)]
Take RS(7,3) over GF(8) with $p(x)=x^3+x+1$ and $b=1$: $n-k=4$, so $t=2$. The generator is
\[
g(x)=(x+\alpha)(x+\alpha^2)(x+\alpha^3)(x+\alpha^4)=x^4+\alpha^3x^3+x^2+\alpha x+\alpha^3,
\]
as can be checked with Table~\ref{tab:ch14:gf8}: $(x+\alpha)(x+\alpha^2)=x^2+\alpha^4x+\alpha^3$,
$(x+\alpha^3)(x+\alpha^4)=x^2+\alpha^6x+1$, and multiplying out gives the coefficients above
(for instance, $\alpha^6+\alpha^4=101\oplus110=011=\alpha^3$).

\emph{Encoding.} The message $(1,\alpha^2,\alpha^5)$ becomes, after systematic encoding, the
codeword $(1,\alpha^2,\alpha^5,\alpha,\alpha,1,\alpha^5)$, the coefficients of $x^6$ down to
$x^0$.

\emph{Channel.} Two symbol errors occur: $\alpha^3$ is added at position $x^5$ and $\alpha^6$ at
position $x^1$, so the received word is $(1,\alpha^5,\alpha^5,\alpha,\alpha,\alpha^2,\alpha^5)$.

\emph{Syndromes.} Evaluating at $\alpha,\alpha^2,\alpha^3,\alpha^4$ gives
$S_1=\alpha^3$, $S_2=\alpha^5$, $S_3=\alpha$, $S_4=\alpha^5$. (Check one from the errors alone:
$S_1=\alpha^3\cdot\alpha^5+\alpha^6\cdot\alpha=\alpha^8+\alpha^7=\alpha+1=\alpha^3$.)

\emph{Error locator.} With $\nu=2$, the recurrence~\eqref{eq:ch14:newton} gives two linear equations,
$S_1\Lambda_2+S_2\Lambda_1=S_3$ and $S_2\Lambda_2+S_3\Lambda_1=S_4$. The determinant is
$S_1S_3+S_2^2=\alpha^4+\alpha^3=\alpha^6$, and Cramer's rule gives
$\Lambda_2=(S_3^2+S_2S_4)/\alpha^6=(\alpha^2+\alpha^3)/\alpha^6=\alpha^5/\alpha^6=\alpha^6$ and
$\Lambda_1=(S_1S_4+S_2S_3)/\alpha^6=(\alpha+\alpha^6)/\alpha^6=\alpha^5/\alpha^6=\alpha^6$. So
$\Lambda(x)=1+\alpha^6x+\alpha^6x^2$, which is indeed $(1+\alpha^5x)(1+\alpha x)$: locators
$X_1=\alpha^5$ and $X_2=\alpha$, positions 5 and 1.

\emph{Error values.} $\Omega(x)=S(x)\Lambda(x)\bmod x^4=\alpha^3+\alpha^3x$ and
$\Lambda'(x)=\alpha^6$. With $b=1$, Forney's formula is $Y=\Omega(X^{-1})/\Lambda'(X^{-1})$:
$Y_1=(\alpha^3+\alpha^3\alpha^2)/\alpha^6=\alpha^2/\alpha^6=\alpha^3$ and
$Y_2=(\alpha^3+\alpha^3\alpha^6)/\alpha^6=\alpha^5/\alpha^6=\alpha^6$. Both errors are found with
their correct values and the message $(1,\alpha^2,\alpha^5)$ is recovered. The same numbers come out
of \texttt{gf.ReedSolomon(7, 3, gf.GF(3), fcr=1)}.
\end{worked}

\subsection{Erasures, and decoding beyond the classical radius}

When the receiver knows that certain symbols are unreliable---a C1 decoder that failed on a CD, a
missing packet, a disk that has died---it can mark them as erasures. The decoder multiplies the
syndrome polynomial by the known \emph{erasure locator} $\Gamma(x)=\prod(1-Z_jx)$ and runs
Berlekamp--Massey on what remains, correcting $\nu$ errors and $e$ erasures whenever
$2\nu+e\le n-k$ (eq.~\ref{eq:ch14:errera}). With erasures only, any $n-k$ lost symbols can be
recovered, the MDS property at work.

Hard-decision decoding stops at $t$ errors. Two directions go further. \emph{List decoding} (Sudan,
1997; Guruswami and Sudan, 1999) interpolates a bivariate polynomial through the received points
and can correct up to $n-\sqrt{nk}$ errors, returning a short list of candidates. The
Koetter--Vardy algorithm (2003) feeds soft reliability information into the same interpolation.
Both gain a few tenths of a decibel to about a decibel in practice, and both remain more of a
research tool than a product feature.

\begin{history}[Five pages that went everywhere]
Reed and Solomon's paper, ``Polynomial codes over certain finite fields,'' occupies five pages of
the \emph{Journal of the Society for Industrial and Applied Mathematics} (1960). It defined the codes
and proved their distance, but its decoding method was impractical for anything but small codes. For
eight years they were a mathematical curiosity. The breakthrough came with Berlekamp's 1968 book
\emph{Algebraic Coding Theory}, whose decoding algorithm made large RS codes practical, and with
Massey's 1969 interpretation of it as shift-register synthesis. By the mid-1980s an RS(255,223) outer code was protecting Voyager's images of the outer planets, and in 1980 Philips and Sony settled on a pair of
shortened RS codes for the compact disc. Reed had already given his name to the Reed--Muller codes in 1954; he went on to a long career
at the University of Southern California.
\end{history}

\mdcfig{ch14_rs_image}{A 96-row image sent as 96 RS(255,223) codewords (one per row), over a
channel that replaces three runs of 260 to 330 consecutive bytes with garbage, about 3.6\% of the
symbols. (a) and (b): codewords sent one after another; each burst wipes out more than 16 symbols of
one or two codewords, the decoder detects that it cannot correct them, and the damage remains. (c)
and (d): the same codewords sent column by column through a block interleaver; each burst is spread
over all 96 codewords, no codeword sees more than 10 errors, and every one is corrected.}

\subsection{Reed--Solomon codes everywhere}

\begin{inpractice}[The compact disc: CIRC]
A CD carries 16-bit stereo audio at 44.1\,kHz. Every six stereo sample pairs (24 bytes) form a
frame. The \term{cross-interleaved Reed--Solomon code} (CIRC) encodes each frame first with a
(28,24) RS code (C2), spreads the 28 bytes over up to 109 successive frames with delays in multiples of four
frames, encodes the result with a (32,28) RS code (C1), and delays alternate bytes by one frame. Both
are shortened from RS(255,251) over GF(256), with $d=5$. The player's C1 decoder corrects single
errors and flags anything worse as erasures; the C2 decoder, seeing the flagged bytes scattered
across many of its codewords, fills in up to four erasures per codeword. The result corrects bursts of
roughly 3500 to 4000 bits, about 2.5\,mm of track, and when even that fails, the player \emph{conceals}
the error by interpolating between neighbouring samples, which hides damage up to roughly 8\,mm
(published figures vary slightly with the assumptions).
Try drilling a small hole in a scratched test disc: the music plays on. The DVD replaced CIRC with a
Reed--Solomon \emph{product code}, (208,192) columns by (182,172) rows, and Blu-ray uses a long
(248,216) RS code interleaved with a ``burst indicator subcode''.
\end{inpractice}

\begin{inpractice}[From QR codes to RAID-6]
A \emph{QR code} (Denso Wave, 1994) stores its data in Reed--Solomon codewords over GF(256) with
the \texttt{0x11D} polynomial. The four error-correction levels, L, M, Q and H, restore roughly 7, 15,
25 and 30\% of damaged codewords, which is why a QR code with a logo pasted over its centre still
scans. \emph{RAID-6} protects an array of $n$ data disks against the loss of any two by storing two
extra blocks, $P=\sum_iD_i$ (plain XOR) and $Q=\sum_i g^iD_i$ computed byte by byte in GF(256) with
$g=\alpha$: exactly the two check symbols of a Reed--Solomon-like MDS code with erasures at known
positions. The Linux kernel's implementation follows H.~Peter Anvin's note ``The mathematics of
RAID-6''. Distributed storage systems scale the idea up: Hadoop's HDFS offers a RS(6,3) erasure
policy (six data blocks, three parity) and large object stores use wider codes, cutting the
storage overhead of triple replication from 200\% to 50\% or less. Digital television uses RS(204,188)
in DVB-S and DVB-T and RS(207,187) in ATSC; optical transport networks (ITU-T G.709) use RS(255,239);
and the IEEE 802.3 Ethernet PHYs of 100\,Gb/s and above use RS(528,514) and RS(544,514) over
GF($2^{10}$), the ``KR4'' and ``KP4'' FECs of Chapter~\ref{ch:baseband}.
\end{inpractice}

%======================================================================
\section{Interleaving and concatenated codes}
\label{sec:ch14:concat}
%======================================================================

Real channels do not make independent errors. A fade on a mobile link, an impulse of noise on a
telephone line, a scratch on a disc or a burst leaving a Viterbi decoder corrupts many
consecutive symbols. A code designed for random errors, faced with a burst longer than it can
correct, simply fails, even though the \emph{average} error rate is well within its power.
Figure~\ref{fig:ch14_rs_image} showed the cure: rearrange the transmitted symbols so that the
burst is shared among many codewords.

\subsection{Block interleavers}

A \term{block interleaver} writes $I$ codewords of length $n$ into the rows of an $I\times n$
array and transmits the array column by column; the de-interleaver does the reverse. A burst of
$b$ consecutive channel symbols then places at most $\lceil b/I\rceil$ errors in any codeword
(Figure~\ref{fig:ch14_interleaving}). The number of rows $I$ is the \term{interleaving depth}. If
each codeword corrects $t$ errors, the interleaved system corrects any single burst of length up
to $tI$. The costs are memory ($In$ symbols at each end) and latency ($2In$ symbol periods end
to end, since the de-interleaver cannot release the first codeword until the whole array has
arrived). CCSDS deep-space telemetry interleaves RS(255,223) codewords to depth 1 to 5 or 8,
correcting single bursts of up to 16 bytes times the depth.

\mdcfig[0.92]{ch14_interleaving}{A burst of nine consecutive channel symbols. Without
interleaving (left) it lands in two codewords, giving four and five errors: too many for a
double-error-correcting code. With a $6\times12$ block interleaver (right), consecutive channel
symbols belong to different codewords, and the same burst leaves at most two errors in any
codeword.}

\subsection{Convolutional interleavers}

A \term{convolutional interleaver}, due to Ramsey and to Forney around 1970--71, achieves the same
spreading with half the memory and delay. The symbol stream is distributed cyclically over $I$
branches with delays $0,M,2M,\dots,(I-1)M$ symbols; the de-interleaver uses the complementary
delays $(I-1)M,\dots,M,0$, so every symbol experiences the same total delay $(I-1)MI$. Adjacent
channel symbols come from branches whose delays differ by $M$, so they are separated by $IM$
positions after de-interleaving. DVB-S and DVB-T use $I=12$ branches with $M=17$ bytes between the
RS(204,188) outer code and the convolutional inner code; $IM=204$, so a burst of up to twelve bytes
puts at most one error into any RS codeword, and bursts of up to $8\times12=96$ bytes are
corrected. ATSC uses a 52-branch interleaver with the same structure.

\begin{pitfall}[Interleavers cost latency, and fading is slow]
An interleaver helps only if its span is long compared with the error bursts. On a mobile channel
at walking speed, a fade can last tens of milliseconds, and an interleaver long enough to break it
up adds more delay than a voice call can tolerate. GSM interleaved speech frames over eight bursts
(about 37\,ms), which works at vehicle speeds but much less well for a pedestrian. Modern systems
attack the problem in frequency instead: OFDM spreads each codeword across many subcarriers that
fade independently (Chapter~\ref{ch:ofdm}), so a single OFDM symbol provides the diversity that a
long time interleaver would otherwise have to buy with delay.
\end{pitfall}

\subsection{Concatenated codes}

In his 1965 doctoral thesis, published as a monograph in 1966, David Forney asked how to build very
long, powerful codes whose decoding complexity grows only polynomially with length. His answer was
to \term{concatenate} two codes (Figure~\ref{fig:ch14_concat_chain}). An \emph{inner} code, typically
a convolutional code with soft-decision Viterbi decoding, faces the raw channel and turns it into a
cleaner ``super-channel''. An \emph{outer} code, typically Reed--Solomon, corrects the residual
errors of the inner decoder. The division of labour is natural: the inner decoder exploits the
soft information that only it can see; its output errors come in bursts, which the byte-oriented
outer code, helped by an interleaver, handles efficiently. Forney proved that concatenated codes
can achieve error probabilities that decrease exponentially with block length at any rate below
capacity, with complexity growing only polynomially.

\begin{figure}[htbp]\centering
\begin{tikzpicture}[node distance=0.3cm and 0.32cm,
  blk/.style={draw=navy,fill=navy!6,thick,minimum height=0.95cm,minimum width=1.45cm,align=center,font=\sffamily\scriptsize},
  outb/.style={draw=accent,fill=accent!6,thick,minimum height=0.95cm,minimum width=1.45cm,align=center,font=\sffamily\scriptsize},
  arr/.style={-{Stealth[length=2mm]},thick,navy}]
\node[outb] (rse) {RS(255,223)\\encoder\\(outer)};
\node[blk,right=of rse] (il) {Interleaver\\depth $I$};
\node[blk,right=of il] (cve) {Conv.\ encoder\\$K=7$, $R=\frac12$\\(inner)};
\node[blk,right=of cve,fill=gray!10,draw=gray] (ch) {BPSK +\\channel};
\node[blk,below=0.9cm of ch] (vit) {Viterbi\\decoder\\(soft)};
\node[blk,left=of vit] (dil) {De-\\interleaver};
\node[outb,left=of dil] (rsd) {RS decoder\\(outer)};
\node[blk,left=of rsd,draw=none,fill=none,minimum width=0.6cm] (sink) {data};
\node[blk,left=of rse,draw=none,fill=none,minimum width=0.6cm] (src) {data};
\draw[arr] (src) -- (rse); \draw[arr] (rse) -- (il); \draw[arr] (il) -- (cve); \draw[arr] (cve) -- (ch);
\draw[arr] (ch) -- node[right,font=\scriptsize]{soft samples} (vit);
\draw[arr] (vit) -- (dil); \draw[arr] (dil) -- (rsd);
\draw[arr] (rsd) -- (sink);
\begin{scope}[on background layer]
\node[draw=navy!50,dashed,rounded corners,fit=(cve)(ch)(vit),inner sep=4pt,label={[font=\sffamily\scriptsize,text=navy]above:inner ``super-channel''}] {};
\end{scope}
\end{tikzpicture}
\caption{The classic concatenated system of CCSDS deep-space telemetry and DVB-S. The inner
convolutional code and soft Viterbi decoder turn a noisy Gaussian channel into a byte channel with
a low error rate; the interleaver breaks the Viterbi decoder's error bursts apart, and the
Reed--Solomon outer code removes what remains.}
\label{fig:ch14_concat_chain}
\end{figure}

Figure~\ref{fig:ch14_concat_ber} quantifies the result for the CCSDS combination. The left panel is a
simulation of the inner decoder: at an inner $E_b/N_0$ of 2\,dB, the Viterbi decoder leaves about one
bit in 200 in error; because those errors cluster, only about one byte in 75 is wrong, far fewer than the
one in 25 that independent bit errors would produce. The right panel feeds
that byte error rate, assuming ideal interleaving, into the RS(255,223) performance
formula. The output error rate plunges from $10^{-2}$ to below $10^{-10}$ over about one decibel,
reaching $10^{-5}$ at about 2.3\,dB overall, including the 0.58\,dB rate loss of the outer code.
That was the state of the art from the late 1970s until turbo codes arrived in 1993, about 2.4\,dB
from the binary-input Shannon limit for the overall rate of $0.437$.

\mdcfig{ch14_concat_ber}{Concatenated coding. Left: simulated bit and byte error rates at the
output of the $K=7$ soft-decision Viterbi decoder. Right: overall bit error rate of the RS(255,223)
plus $K=7$ system with ideal interleaving (outer-code performance computed from the inner byte error
rate), compared with the inner code alone.}

\begin{history}[From Mariner to Voyager]
Deep space was the proving ground of coding because every decibel there is worth a fortune: it can
be bought only with a bigger antenna on Earth or a bigger transmitter on the spacecraft. Mariner 6
and 7 (1969) and Mariner 9 (1971), which mapped Mars from orbit, sent their images with the
(32,6) first-order Reed--Muller code, each 6-bit pixel expanded to 32 bits and decoded on the ground
by a fast Hadamard transform in a machine called the ``Green machine'' after its designer, R.~R.~Green of JPL.
Pioneer 10 and 11 used long convolutional codes with sequential decoding. Voyager (1977) used the
$K=7$, rate-1/2 convolutional code with Viterbi decoding, and for the encounters with Uranus (1986) and
Neptune (1989), where the signal was weakest, relied on concatenation with the RS(255,223) outer
code together with image data compression. The same concatenated system was standardised by the
Consultative Committee for Space Data Systems (CCSDS) and flew on missions for decades; it survives
in the CCSDS recommendation (CCSDS 131.0-B) alongside the turbo and LDPC codes that superseded it.
\end{history}

\subsection{Product codes}

Peter Elias proposed an even simpler way to build a long code from short ones in 1954. Arrange
$k_1k_2$ information bits in a $k_1\times k_2$ array, encode each row with an $(n_2,k_2,d_2)$ code
and then each column (including the parity columns) with an $(n_1,k_1,d_1)$ code. The result is an
$(n_1n_2,k_1k_2,d_1d_2)$ \term{product code}. The minimum distance multiplies, while decoding can
still be done row by row and column by column, iterating between the two. With hard-decision row
and column decoders a product code already performs well; with soft-input soft-output component
decoders exchanging reliability information (Pyndiah's ``turbo product codes'' of 1994--98) it
comes within about a decibel of capacity. Product codes protect DVD sectors, and their descendants,
staircase codes and related ``spatially coupled'' product-like codes, carry 100 to 400\,Gb/s optical
links (Chapter~\ref{ch:wireline}).

\subsection{Reed--Muller codes}

The Reed--Muller codes, found by David Muller in 1954 with a majority-logic decoder by Irving Reed
in the same year, are among the oldest and most beautiful families. The code $\mathrm{RM}(r,m)$ has
length $n=2^m$, dimension $k=\sum_{i=0}^{r}\binom mi$ and minimum distance $2^{m-r}$. Its codewords are
the truth tables of Boolean functions of $m$ variables of degree at most $r$. The first-order codes
$\mathrm{RM}(1,m)$ are the \emph{biorthogonal} codes: $\mathrm{RM}(1,5)$ is the (32,6,16) Mariner code, whose
64 codewords are the rows of a $32\times32$ Hadamard matrix and their complements, so soft
maximum-likelihood decoding is a single fast Hadamard transform. The recursive construction
\begin{equation}
\mathrm{RM}(r,m)=\big\{(\vect{u},\vect{u}+\vect{v}):\ \vect{u}\in\mathrm{RM}(r,m-1),\ \vect{v}\in\mathrm{RM}(r-1,m-1)\big\}
\label{eq:ch14:uuv}
\end{equation}
is the same $(\vect{u},\vect{u}+\vect{v})$ structure that underlies polar codes: both codes select rows of
the Kronecker power of $\left[\begin{smallmatrix}1&0\\1&1\end{smallmatrix}\right]$, Reed--Muller by
row weight and polar by channel reliability (Chapter~\ref{ch:moderncodes}). In a pleasing twist, the
Reed--Muller codes themselves were shown to achieve capacity on the erasure channel in 2016--17, and a
proof for all binary memoryless symmetric channels followed in 2023. LTE and NR use a (32,$O$)
Reed--Muller-based code for short uplink control reports.

%======================================================================
\section{Convolutional codes}
\label{sec:ch14:conv}
%======================================================================

Block codes chop the data into independent blocks. Peter Elias's convolutional codes of 1955
take the opposite view: the encoder is a small finite-state machine that slides along an
unbounded data stream, and every output bit depends on the current input and a few previous ones.
There are no block boundaries, the encoder is a handful of flip-flops and XOR gates, and, crucially,
the codes can be decoded with soft decisions by a simple, regular algorithm.

\subsection{The encoder}

A rate-$1/n$ \term{convolutional encoder} contains a shift register holding the last $K-1$ input
bits. For each input bit $u_k$ it emits $n$ output bits, each the modulo-2 sum of a fixed subset of
$u_k,u_{k-1},\dots,u_{k-K+1}$. $K$ is the \term{constraint length} (some authors use the memory
$K-1$ instead: read every paper's definition). Figure~\ref{fig:ch14_conv_encoder} shows the
four-state, rate-1/2 code with $K=3$ that serves as the textbook example:
\begin{equation}
c_k^{(1)}=u_k\oplus u_{k-1}\oplus u_{k-2},\qquad c_k^{(2)}=u_k\oplus u_{k-2}.
\label{eq:ch14:conv75}
\end{equation}
Each output stream is the input convolved (modulo 2) with a \term{generator sequence}, here
$\vect{g}^{(1)}=(1,1,1)$ and $\vect{g}^{(2)}=(1,0,1)$, whence the name. Generators are written in
octal, $(7,5)$ for this code, or as polynomials in the delay operator $D$:
$g^{(1)}(D)=1+D+D^2$ and $g^{(2)}(D)=1+D^2$. In the $D$-transform domain encoding is multiplication,
$c^{(j)}(D)=u(D)g^{(j)}(D)$, and a convolutional code is a linear code over the field of rational
functions in $D$, a view that makes much of the theory parallel that of block codes.

\begin{figure}[htbp]\centering
\begin{tikzpicture}[>=Stealth,
  reg/.style={draw=navy,thick,fill=navy!6,minimum width=1.0cm,minimum height=0.8cm,font=\sffamily\small},
  xr/.style={draw=accent,thick,circle,inner sep=0pt,minimum size=0.5cm,font=\small},
  lab/.style={font=\sffamily\scriptsize,text=navy}]
\coordinate (in) at (0,0);
\node[reg] (s1) at (2.2,0) {$u_{k-1}$};
\node[reg] (s2) at (4.0,0) {$u_{k-2}$};
\coordinate (t0) at (1.0,0);
\node[xr] (xa) at (2.2,1.6) {$+$};
\node[xr] (xb) at (2.2,-1.6) {$+$};
\draw[->,thick,navy] (in) node[left,lab]{$u_k$} -- (s1);
\draw[->,thick,navy] (s1) -- (s2);
\fill[navy] (t0) circle (1.5pt);
\draw[->,thick,navy] (t0) |- (xa);
\draw[->,thick,navy] (t0) |- (xb);
\draw[->,thick,navy] (s1) -- (xa);
\draw[thick,navy] (s2.east) -- ++(0.5,0) coordinate (t2);
\fill[navy] (t2) circle (1.5pt);
\draw[->,thick,navy] (t2) |- (xa);
\draw[->,thick,navy] (t2) |- (xb);
\draw[->,thick,accent] (xa) -- ++(3.4,0) node[right,lab,text=accent]{$c^{(1)}_k$, $\vect g^{(1)}=111=7_8$};
\draw[->,thick,accent] (xb) -- ++(3.4,0) node[right,lab,text=accent]{$c^{(2)}_k$, $\vect g^{(2)}=101=5_8$};
\node[lab,align=left,anchor=west] at (5.2,0) {state $=(u_{k-1},u_{k-2})$};
\end{tikzpicture}
\caption{The rate-1/2, constraint-length-3 convolutional encoder with generators $(7,5)$ in octal.
The two register stages hold the encoder state; each input bit produces two output bits, which are
multiplexed onto the channel.}
\label{fig:ch14_conv_encoder}
\end{figure}

\subsection{State diagram and trellis}

The encoder's state is the register contents $(u_{k-1},u_{k-2})$, four states in all; each input
bit moves it to a new state and produces an output pair. Figure~\ref{fig:ch14_state} draws this as
a \term{state diagram}. Unrolling the state diagram in time gives the \term{trellis} of
Figure~\ref{fig:ch14_viterbi_example}: one column of states per time step, with each state joined to
the two states it can reach. Every path through the trellis is a codeword, and every codeword is a
path. The trellis is the central object of the decoding theory; Forney introduced the name in 1973.

\begin{figure}[htbp]\centering
\begin{tikzpicture}[>=Stealth,
  st/.style={draw=navy,thick,circle,fill=navy!6,minimum size=0.95cm,font=\sffamily\small},
  lab/.style={font=\sffamily\scriptsize,fill=white,inner sep=1pt}]
\node[st] (a) at (0,0) {00};
\node[st] (b) at (2.6,1.5) {10};
\node[st] (c) at (2.6,-1.5) {01};
\node[st] (d) at (5.2,0) {11};
\draw[->,thick,navy] (a) edge[loop left,looseness=6] node[lab,left]{0/00} (a);
\draw[->,thick,accent,dashed] (a) -- node[lab,above left]{1/11} (b);
\draw[->,thick,navy] (b) to[bend left=18] node[lab,right,pos=0.4]{0/10} (c);
\draw[->,thick,accent,dashed] (b) -- node[lab,above right]{1/01} (d);
\draw[->,thick,navy] (c) -- node[lab,below left]{0/11} (a);
\draw[->,thick,accent,dashed] (c) to[bend left=18] node[lab,left,pos=0.4]{1/00} (b);
\draw[->,thick,navy] (d) -- node[lab,below right]{0/01} (c);
\draw[->,thick,accent,dashed] (d) edge[loop right,looseness=6] node[lab,right]{1/10} (d);
\end{tikzpicture}
\caption{State diagram of the $(7,5)$ code. States are $(u_{k-1},u_{k-2})$; each branch is labelled
input/output. Solid branches are input 0, dashed branches input 1.}
\label{fig:ch14_state}
\end{figure}

\subsection{Free distance and the transfer function}

Because the code is linear, its distance properties are those of the all-zero path. An
\emph{error event} is a path that leaves the all-zero state and later returns to it; the smallest
output weight of any error event is the \term{free distance} $d_{\text{free}}$, which plays the role
of $d_{\min}$. For the $(7,5)$ code, the path $00\to10\to01\to00$ (input $100$) produces outputs
$11,10,11$, weight 5, and no shorter or lighter detour exists, so $d_{\text{free}}=5$.

The complete distance structure is captured by the \term{transfer function}. Label each branch of
the state diagram with $D^{w}N^{i}$, where $w$ is its output weight and $i$ its input weight, split
the zero state into a source and a sink, and solve the flow-graph equations. Writing $X_s$ for the
state variables of the $(7,5)$ code,
\begin{gather*}
X_{10}=D^2N\,X_{\text{in}}+N\,X_{01},\qquad X_{01}=D\,X_{10}+D\,X_{11},\\
X_{11}=DN\,X_{10}+DN\,X_{11},\qquad X_{\text{out}}=D^2X_{01},
\end{gather*}
and eliminating gives
\begin{equation}
T(D,N)=\frac{X_{\text{out}}}{X_{\text{in}}}=\frac{D^5N}{1-2DN}=D^5N+2D^6N^2+4D^7N^3+\cdots .
\label{eq:ch14:tf}
\end{equation}
There is one error event of weight 5 (with one input error), two of weight 6 (each with two), four of
weight 7, and so on. The derivative $\partial T/\partial N$ at $N=1$,
\[
\frac{D^5}{(1-2D)^2}=D^5+4D^6+12D^7+32D^8+\cdots=\sum_dB_dD^d ,
\]
gives the total number $B_d$ of information-bit errors on all weight-$d$ events, which is exactly
what a bit-error-rate bound needs. For large codes the transfer function is computed numerically by
a search over the trellis, as \lab{book/figscripts/ch14\_figs.py} does; Table~\ref{tab:ch14:conv}
lists the results for the best-known rate-1/2 codes.

\begin{table}[htbp]\centering\small
\caption{The best rate-1/2 convolutional codes of small constraint length (maximum free distance,
found by computer search in the 1970s). $B_d$ counts information-bit errors on weight-$d$ error
events; the asymptotic coding gain is $10\log_{10}(Rd_{\text{free}})$.}
\label{tab:ch14:conv}
\begin{tabular}{@{}c c c c l c@{}}
\toprule
$K$ & states & generators (octal) & $d_{\text{free}}$ & $B_d$ for $d=d_{\text{free}},d_{\text{free}}+1,\dots$ & gain (dB)\\
\midrule
3 & 4 & (7, 5) & 5 & 1, 4, 12, 32, \dots & 3.98\\
5 & 16 & (23, 35) & 7 & 4, 12, 20, 72, \dots & 5.44\\
7 & 64 & (171, 133) & 10 & 36, 0, 211, 0, 1404, 0, 11633, \dots & 6.99\\
9 & 256 & (561, 753) & 12 & 33, 0, 281, 0, 2179, 0, 15035, \dots & 7.78\\
\bottomrule
\end{tabular}
\end{table}

\subsection{Catastrophic codes and other properties}

Not every choice of generators is a good code. If the generators share a common polynomial factor
other than a power of $D$, the encoder is \term{catastrophic}: some input sequence of infinite weight
produces an output of finite weight, so a finite number of channel errors can cause an infinite
number of decoded errors. The code $(3,5)_8$, with $g^{(1)}=1+D$ and $g^{(2)}=1+D^2=(1+D)^2$, is the
standard example: the all-ones input produces the output $11,10,00,00,\dots$, of weight three.
Massey and Sain showed in 1968 that a rate-$1/n$ encoder is non-catastrophic exactly when the
greatest common divisor of its generators is a power of $D$. Code-search tables list only
non-catastrophic codes.

Convolutional encoders can also be \emph{systematic} (one output equal to the input) and
\emph{recursive} (with feedback, like the CRC division circuit). Non-recursive systematic codes have
noticeably smaller free distance for the same $K$, which is why the classical standards use
non-systematic feed-forward encoders. Recursive systematic encoders with the same free distance as the
best feed-forward ones exist, and their infinite impulse response turns out to be essential for turbo
codes (Chapter~\ref{ch:moderncodes}).

\subsection{The industry-standard $K=7$ code}

One code dominates the history of convolutional coding: the rate-1/2, $K=7$ code with generators
$(171,133)_8$, found by Odenwalder in 1970. With 64 states it was small enough to decode in 1970s
hardware, and with $d_{\text{free}}=10$ it gives about 5\,dB of coding gain. NASA adopted it for
Voyager; CCSDS standardised it for space telemetry; it is the inner code of DVB-S and DVB-T, and the
convolutional code of IEEE 802.11a/g/n/ac Wi-Fi, where it is still the mandatory code; and it
appears, with a third generator (165) and rate 1/3, as the tail-biting control-channel code of LTE.
Figure~\ref{fig:ch14_conv_ber} compares it with the other codes of Table~\ref{tab:ch14:conv}. Each
increase of two in $K$ buys roughly 0.5 to 1\,dB, at four times the decoding complexity.

\mdcfig[0.9]{ch14_conv_ber}{Simulated bit error rates of the rate-1/2 codes of
Table~\ref{tab:ch14:conv} with soft-decision Viterbi decoding (points), and their union
bounds~\eqref{eq:ch14:ub} (dotted), which become tight below about $10^{-5}$.}

\begin{pitfall}[Which way round is (171,133)?]
Octal generators hide two conventions. The first is the order of the outputs: 802.11 lists
$g_0=133$, $g_1=171$ and sends the 133 output first; DVB and CCSDS list $G_1=171$, $G_2=133$ and
send 171 first. The second is the bit order within the generator: with the most significant bit
tapping the \emph{current} input (the convention of this book), $171_8=1111001$; with the least significant bit
tapping the current input, the same taps read $1001111=117_8$, and a paper using that convention
will call this code $(117,155)$. CCSDS additionally inverts the second output. Any one of these
mismatches produces a decoder that fails completely, and the quickest diagnosis is to encode an
impulse (a single one followed by zeros) and compare the output with the standard's.
\end{pitfall}

\subsection{Puncturing}

A rate-1/2 code is often too slow. Designing new rate-2/3 or 3/4 codes with their own trellises is
possible but unnecessary: \term{puncturing} deletes some coded bits in a periodic pattern, raising the
rate while keeping the decoder of the mother code. The Viterbi decoder simply inserts neutral values
(zero log-likelihood ratios, or erasures) in the deleted positions and runs as usual. IEEE 802.11a
obtains rate 2/3 by deleting every second bit of the $g_1$ output, and rate 3/4 by sending
$A_1B_1A_2B_3$ out of every six bits $A_1B_1A_2B_2A_3B_3$; 802.11n adds rate 5/6, and DVB-S uses
rates 1/2, 2/3, 3/4, 5/6 and 7/8. Figure~\ref{fig:ch14_puncturing} shows the price: each step up in
rate costs between 0.6 and 1\,dB, and the rate-5/6 code still gains about 3\,dB over uncoded transmission at
$10^{-5}$ while carrying 83\% as much data.

\mdcfig[0.88]{ch14_puncturing}{The $K=7$ mother code punctured to the rates of IEEE 802.11a/n,
simulated with soft-decision Viterbi decoding. Energy is counted per information bit, so the curves
include the rate change.}

Hagenauer's rate-compatible punctured convolutional (RCPC) codes of 1988 go one step further: a family
of puncturing patterns chosen so that the bits kept at each rate are a subset of those kept at every
lower rate. A transmitter can then send the highest-rate version first and, if the receiver fails,
send only the additional bits needed for the next lower rate. That is incremental redundancy, the
idea behind HARQ (Section~\ref{sec:ch14:arq}).

\subsection{Termination and tail-biting}

A convolutional encoder is normally \emph{terminated} by appending $K-1$ zeros to each block, returning
it to the zero state so that the decoder knows where the path ends. For a 1500-byte Wi-Fi packet the six
tail bits are negligible, but for a 40-bit control message they cost 15\% of the rate. A
\term{tail-biting} encoder avoids the overhead by initialising its register with the \emph{last} $K-1$
bits of the block, so the path starts and ends in the same, unknown, state. The decoder handles the
circular trellis by running the Viterbi algorithm around it more than once, or by trying the likely
starting states. LTE's control and broadcast channels use tail-biting for exactly this reason.

%======================================================================
\section{The Viterbi algorithm}
\label{sec:ch14:viterbi}
%======================================================================

Andrew Viterbi published his algorithm in 1967 as a device for proving error bounds for convolutional
codes; within a few years it was recognised as the \emph{maximum-likelihood} decoder, efficient enough to
build, and it became one of the most widely executed algorithms in history. Its reach extends far
beyond coding: the same dynamic program decodes intersymbol interference (Chapter~\ref{ch:equalization}),
detects partial-response signals in disk drives (Chapter~\ref{ch:baseband}), aligns DNA sequences, tags
parts of speech and recognises speech with hidden Markov models.

\subsection{Maximum-likelihood decoding as a shortest path}

The decoder receives a sequence $\vect{y}$ and must find the codeword, that is the trellis path, that
maximises $p(\vect{y}\mid\vect{c})$. For a memoryless channel the likelihood factors over time, so its
logarithm is a sum of \term{branch metrics}, one per trellis branch. For the binary symmetric channel
the branch metric can be the Hamming distance between the branch's output bits and the received bits;
for BPSK in Gaussian noise it can be the squared Euclidean distance, or, dropping terms common to all
branches, minus the correlation $\sum_jx_jy_j$. Maximum likelihood becomes a shortest-path problem on the
trellis: find the path of minimum total metric.

There are $2^{k}$ paths through a trellis of length $k$, but the search is easy because of a
simple observation. Suppose two paths arrive at the same state at the same time. Whatever happens later,
both can be extended by exactly the same continuations, which add exactly the same metrics. So the path
with the larger metric so far can never become the best, and can be discarded. At each time and each
state only one \term{survivor} needs to be kept. The algorithm is therefore:
\begin{enumerate}
\item \textbf{Add.} For each state, add the branch metric of each incoming branch to the path metric of
the state it comes from.
\item \textbf{Compare.} Compare the candidate metrics arriving at the state.
\item \textbf{Select.} Keep the smallest as the new path metric and record which branch won (the
\emph{decision bit}).
\item Repeat for every time step; at the end, start from the best (or known final) state and
\textbf{trace back} through the stored decisions to recover the input sequence.
\end{enumerate}
The add--compare--select (\term{ACS}) operation is the heart of every Viterbi decoder. For a code with
$2^{K-1}$ states it runs $2^{K-1}$ times per decoded bit: 64 ACS operations per bit for the $K=7$ code.
The work grows linearly with the length of the data and exponentially with the constraint length, which
is why practical Viterbi decoders stop at $K\approx9$ (256 states) and why JPL's experimental ``Big
Viterbi Decoder'' for $K=15$, with 16\,384 states, was a heroic exception.

\begin{worked}[Viterbi decoding by hand]
The $(7,5)$ encoder of Figure~\ref{fig:ch14_conv_encoder}, started in state 00, encodes $u=1011$
followed by two zero tail bits as $11\;10\;00\;01\;01\;11$. The channel flips two bits, and the decoder
receives $11\;\mathbf{11}\;00\;01\;\mathbf{11}\;11$. We decode with Hamming branch metrics
(Figure~\ref{fig:ch14_viterbi_example}).

\emph{Step 1} (received 11). From state 00 (metric 0), input 0 gives output 00 at distance 2, input 1 gives
output 11 at distance 0. Metrics: state 00 $\to2$, state 10 $\to0$.

\emph{Step 2} (received 11). From 00 (2): to 00 via 00, $2+2=4$; to 10 via 11, $2+0=2$. From 10 (0): to 01
via 10, $0+1=1$; to 11 via 01, $0+1=1$. Metrics: 00:4, 10:2, 01:1, 11:1.

\emph{Step 3} (received 00). Now every state has two incoming branches, and the ACS begins. State 00 can be
reached from 00 (4, output 00, $+0$) or from 01 (1, output 11, $+2$): the survivor comes from 01 with
metric 3. State 10: from 00 via 11, $4+2=6$, or from 01 via 00, $1+0=1$; survivor 1. State 01: from 10
via 10, $2+1=3$, or from 11 via 01, $1+1=2$; survivor 2. State 11: from 10 via 01, $2+1=3$, or from 11
via 10, $1+1=2$; survivor 2.

\emph{Steps 4--6} continue in the same way (the figure shows every metric). After the sixth step, state
00 has metric 2: the best codeword differs from the received sequence in exactly two places. Tracing back
from state 00 through the recorded survivors visits $00\to10\to01\to10\to11\to01\to00$, which
corresponds to inputs $1\,0\,1\,1\,0\,0$. Both channel errors are corrected. They were far enough apart
(two branches) that no competing path came within distance two; had they been adjacent, a weight-5 error
event could have won.
\end{worked}

\mdcfig[0.95]{ch14_viterbi_example}{The Viterbi algorithm on the $(7,5)$ trellis for the worked example.
States (rows) are $u_{k-1}u_{k-2}$; solid branches are input 0 and dashed branches input 1. The
number in each node is the accumulated Hamming path metric. Dark branches are survivors; the shaded path
is the traceback from state 00 at the end. Received pairs in red contain a channel error.}

\subsection{Soft decisions, quantisation and metric arithmetic}

With soft decisions the branch metric is the correlation between the branch's BPSK symbols and the
received samples, or equivalently a sum of log-likelihood ratios with signs set by the branch's output
bits; commlib's decoder uses exactly this. No knowledge of the noise variance is needed on the AWGN
channel, because scaling every metric by the same factor does not change the winner. Hardware
quantises the samples to a few bits: the right panel of Figure~\ref{fig:ch14_viterbi_practical} shows that
three bits (eight levels) come within about 0.25\,dB of unquantised decoding and four bits within
about 0.1\,dB, while two bits cost about 0.7\,dB and hard decisions (one bit) more than 2\,dB. Path metrics grow without
bound, but only their differences matter, and the spread of metrics across states is bounded (by
roughly $d_{\text{free}}$ times the maximum branch metric), so hardware either subtracts the minimum
periodically or, more elegantly, lets the metrics wrap around in modular two's-complement arithmetic
and compares them modulo the register size.

\subsection{Traceback and decision depth}

Waiting for the end of a long stream before tracing back is impractical: memory and latency would be
unbounded. Fortunately, survivors \emph{merge}. If one traces back from all states at time $t$, the
paths coincide a few constraint lengths earlier with high probability, so the decision for time $t-D$ can
be released after $D$ steps. A \term{decision depth} (traceback depth) of about five constraint lengths
is the classic rule of thumb for rate-1/2 codes; punctured, higher-rate codes, whose paths take longer to
separate, need more, typically seven to ten constraint lengths. The left panel of
Figure~\ref{fig:ch14_viterbi_practical} shows the effect for the $K=7$ code: short depths are
disastrous, at $D=35=5K$ the bit error rate is within about 30\% of full traceback (a loss of only
a few hundredths of a decibel on these steep curves), and beyond about $D=45$ the difference
vanishes. Hardware implements traceback with a memory of decision
bits read backwards in blocks, or with the register-exchange method, which copies whole survivor
sequences at every step and trades power for latency.

\mdcfig{ch14_viterbi_practical}{Implementation choices for the $K=7$ Viterbi decoder. Left: bit
error rate against decision depth, with full-traceback performance dotted. Right: the effect of
quantising the received samples to 1, 2, 3 and 4 bits.}

\subsection{Performance: the union bound}

The decoder errs when the correct path is beaten by an error event. The probability that a specific
event at distance $d$ beats the correct path is $\Q\big(\sqrt{2dRE_b/N_0}\big)$ with soft decisions, and a
union bound over all events, weighted by the information bits each one corrupts, gives
\begin{equation}
P_b\le\sum_{d\ge d_{\text{free}}}B_d\,\Q\!\left(\sqrt{\frac{2dRE_b}{N_0}}\right),
\label{eq:ch14:ub}
\end{equation}
with $B_d$ from the transfer function (Table~\ref{tab:ch14:conv}). The bound diverges at low
$E_b/N_0$, where the events are far from disjoint, but becomes tight below about $10^{-5}$
(Figure~\ref{fig:ch14_conv_ber}). For hard decisions the pairwise probability is replaced by the
probability that more than half of $d$ binary symmetric channel bits are flipped, often bounded by
the Bhattacharyya form $\big(2\sqrt{p(1-p)}\big)^d$. At high SNR the first term dominates, giving
$P_b\approx B_{d_{\text{free}}}\Q\big(\sqrt{2d_{\text{free}}RE_b/N_0}\big)$ and the asymptotic coding gain
$10\log_{10}(Rd_{\text{free}})$.

\begin{worked}[The coding gain of the $K=7$ code at $10^{-5}$]
Uncoded BPSK needs $\Q(\sqrt{2E_b/N_0})=10^{-5}$, that is $\sqrt{2E_b/N_0}=4.265$ and
$E_b/N_0=9.10$, or 9.59\,dB.

For the $K=7$ code, try $E_b/N_0=4.4$\,dB $=2.754$. With $R=\tfrac12$ the argument of the
$d$-th term is $\sqrt{dE_b/N_0}$. The first four terms of~\eqref{eq:ch14:ub} are
$36\,\Q(\sqrt{27.5})=36\times7.7\times10^{-8}=2.8\times10^{-6}$,
$211\,\Q(\sqrt{33.0})=9.5\times10^{-7}$, $1404\,\Q(\sqrt{38.6})=3.7\times10^{-7}$ and
$11633\,\Q(\sqrt{44.1})=1.8\times10^{-7}$, and the sum of all terms is about $4.4\times10^{-6}$.
The terms fall off slowly, because the multiplicities grow almost as fast as the $\Q$-function
shrinks: the ``nearest neighbours'' penalty. Solving for $10^{-5}$ gives $E_b/N_0\approx4.2$\,dB
from the bound, and the simulation of Figure~\ref{fig:ch14_conv_ber} gives about 4.1\,dB
(practical decoders with three-bit quantisation need 0.2 to 0.3\,dB more). The coding gain is
therefore about $9.6-4.1\approx5.5$\,dB, compared with the asymptotic
$10\log_{10}(0.5\times10)=7.0$\,dB. At $10^{-7}$ the bound gives 5.4\,dB against 11.3\,dB uncoded, a
gain of almost 6\,dB: coding gain grows as the target error rate falls. With hard decisions the gain at
$10^{-5}$ shrinks to about 3\,dB (Figure~\ref{fig:ch14_hard_soft}).
\end{worked}

\begin{history}[The algorithm that was not patented]
Viterbi, born in Bergamo in 1935, emigrated to the United States as a child, studied at MIT and worked
at JPL before joining UCLA. His 1967 paper in the \emph{IEEE Transactions on Information Theory}
presented the algorithm as a proof technique; Jim Omura observed in 1969 that it is dynamic programming,
and Forney showed in 1973, in a celebrated tutorial in the \emph{Proceedings of the IEEE}, that it is the
maximum-likelihood sequence detector for any finite-state process, not just for codes. Viterbi did not
patent it. In 1968 he co-founded Linkabit with Irwin Jacobs and Leonard Kleinrock, and Linkabit built
some of the first practical Viterbi decoders for military satellite links; Heller and Jacobs's 1971
paper, showing that short-constraint-length codes with soft Viterbi decoding gave 5\,dB or more of gain in
practice, persuaded the industry. In 1985 Viterbi and Jacobs co-founded Qualcomm, whose CDMA phones
decoded the $K=9$ convolutional codes of IS-95 with, of course, the Viterbi algorithm.
\end{history}

\begin{inpractice}[Viterbi decoders in silicon]
A 64-state Viterbi decoder for 802.11a at 54\,Mb/s needs about 3.5 billion ACS operations per second,
which a 2000-era chip provided with 64 parallel ACS units at a modest clock. Faster decoders process two
trellis steps per clock (``radix-4'', four candidates per state) or split the stream into overlapping
blocks decoded in parallel; 802.11n/ac chips decode hundreds of megabits per second this way, although
the faster Wi-Fi modes now prefer LDPC. The convolutional code also survives where frames are short and
latency matters: the control channels of LTE, GNSS navigation messages (GPS L2C and L5 use the same
$K=7$ code) and countless proprietary radios. Every software-defined radio toolkit, including GNU Radio,
ships a Viterbi decoder for it.
\end{inpractice}

\subsection{Beyond hard outputs}

The Viterbi algorithm outputs a sequence of bits with no indication of their reliability. The
soft-output Viterbi algorithm (SOVA) of Hagenauer and Hoeher (1989) estimates the reliability of each
decision from the metric difference of the competing path, and the BCJR algorithm of Bahl, Cocke,
Jelinek and Raviv (1974) computes exact a-posteriori probabilities for every bit with a forward and a
backward recursion over the same trellis. BCJR was a curiosity for almost twenty years---it costs roughly
three times as much as Viterbi and gives almost the same bit error rate---until turbo decoding made
soft outputs essential. Chapter~\ref{ch:moderncodes} takes up the story.

%======================================================================
\section{Sequential decoding}
\label{sec:ch14:sequential}
%======================================================================

Before the Viterbi algorithm, and for a while alongside it, the preferred way to decode convolutional
codes was \term{sequential decoding}, proposed by John Wozencraft at MIT in 1957. The idea is to
explore the code \emph{tree} (the trellis without merging) like a hiker following the most promising trail
and backing up when it seems wrong. Because the decoder examines only a few paths when the noise is low,
its average complexity is almost independent of the constraint length, so codes with $K=30$ or more,
far too big for Viterbi decoding, become practical, and their error probabilities are correspondingly tiny.

The decoder ranks partial paths by the \term{Fano metric}, which for each coded bit adds
$\log_2\big[p(y\mid x)/p(y)\big]-R$: the information the received bit gives about the hypothesis, minus
the rate. On the correct path the metric tends to increase; on a wrong path it tends to decrease. The
\emph{stack algorithm} (Zigangirov, 1966; Jelinek, 1969) keeps a sorted list of partial paths and always
extends the best. The \emph{Fano algorithm} (1963) achieves the same with almost no memory, moving
forwards and backwards along a single path against a running threshold.

The catch is that the amount of computation is a random variable. When the noise is bad, the decoder
backs up and searches, and the number of computations per decoded bit has a Pareto (power-law)
distribution. Its mean is finite only for rates below the \term{computational cutoff rate} $R_0$. For
BPSK with soft decisions, $R_0=1-\log_2\!\big(1+e^{-E_s/N_0}\big)$, and at rate 1/2 this requires
$E_b/N_0\approx2.5$\,dB; with hard decisions, $R_0=1-\log_2\big(1+2\sqrt{p(1-p)}\big)$ requires about
4.6\,dB. When the computation exceeds the decoder's buffer, the frame must be erased. For two decades
$R_0$ was regarded as the ``practical capacity'' of a channel, the rate beyond which no practical decoder
could go, and closing the gap between $R_0$ and capacity (about 2.3\,dB at rate 1/2 on the Gaussian
channel) was one of the motivations of concatenated coding. Turbo and LDPC codes finally showed that
$R_0$ is not a fundamental limit at all.

Sequential decoding flew on NASA's Pioneer missions, and sequential decoders were built for military and
satellite modems in the 1960s and 1970s. It lost to Viterbi decoding for three reasons: the variable
computation complicated real-time hardware, Viterbi decoding is naturally soft-decision and parallel,
and a 64-state Viterbi decoder gave nearly as much gain as a long code decoded sequentially at the
error rates that mattered.

%======================================================================
\section{ARQ and hybrid ARQ}
\label{sec:ch14:arq}
%======================================================================

When a return channel exists, the receiver can do something no forward code can: ask again. ARQ
protocols are as old as the telegraph operator who tapped ``repeat'', and they are the reason the
Internet delivers files without a single wrong bit despite running over links with residual error rates
of one in a million or worse.

\subsection{The three classical protocols}

All ARQ schemes send frames protected by a CRC and wait for acknowledgements (ACK) or negative
acknowledgements (NAK). They differ in what the transmitter does while it waits and after an error. Let
each frame take one time unit to send, let the one-way propagation delay be $a$ frame times, and let
each frame be in error with probability $P$.
\begin{description}[style=nextline,leftmargin=1.2em,font=\sffamily\bfseries\small\color{navy}]
\item[Stop-and-wait.] Send a frame, wait for its ACK, and resend it if a NAK or a timeout arrives. Each
attempt occupies $1+2a$ time units and $1/(1-P)$ attempts are needed on average, so the throughput
efficiency is
\begin{equation}
\eta_{\text{SW}}=\frac{1-P}{1+2a}.
\end{equation}
Simple and memory-free, it is hopeless on long links: with $a=10$ the channel is idle 95\% of the time.
\item[Go-back-$N$.] Send frames continuously, up to a window of $N\ge1+2a$ unacknowledged frames. On an
error, resend the failed frame and \emph{everything after it}. The receiver needs no buffer, but each
error wastes a round trip of frames:
\begin{equation}
\eta_{\text{GBN}}=\frac{1-P}{1+2aP}.
\end{equation}
\item[Selective repeat.] Send continuously and resend only the frames that failed. The receiver must
buffer out-of-order frames and reassemble them, but the efficiency is the best possible:
\begin{equation}
\eta_{\text{SR}}=1-P.
\end{equation}
\end{description}
Figure~\ref{fig:ch14_arq} (left) compares them. TCP's cumulative acknowledgements behave like go-back-$N$,
with selective acknowledgements (SACK) adding selective repeat; 802.11 uses stop-and-wait for single frames
and a selective ``block acknowledgement'' for aggregated frames; LTE and NR run a selective-repeat ARQ in
the radio link control layer on top of HARQ.

\subsection{Hybrid ARQ}

Pure ARQ throws away a failed frame. Hybrid ARQ keeps it. In \term{Chase combining} (Chase, 1985) every
retransmission repeats the same codeword and the receiver adds the log-likelihood ratios of all copies
before decoding, so the effective signal-to-noise ratio grows with each attempt. In \term{incremental
redundancy} (IR) each retransmission sends \emph{new} parity bits---a different puncturing of a low-rate
mother code, as in the RCPC codes above---so the receiver decodes an ever lower-rate code. IR gains more,
because a lower-rate code is worth more than a repeated one: two transmissions at signal-to-noise ratio
$\gamma$ give an accumulated mutual information of $2\log_2(1+\gamma)$ with IR, but only $\log_2(1+2\gamma)$
with Chase combining. Figure~\ref{fig:ch14_arq} (right) shows the effect on a Rayleigh block-fading
channel with an initial rate of 3\,b/s/Hz: IR is worth 2 to 4\,dB over Chase combining at low and moderate SNR,
and both are far better than discarding failed frames, even though the transmitter knows nothing about
the instantaneous channel. At high SNR every scheme saturates at the initial rate $R_0$, which is why real
systems pair HARQ with link adaptation that chooses $R_0$ from channel-quality reports. HARQ turns a
fixed-rate code into a rate-adaptive one, driven by the receiver's CRC.

\mdcfig{ch14_arq}{Left: throughput efficiency of the three ARQ protocols with a one-way delay of ten frame
times. Right: throughput of HARQ on a Rayleigh block-fading channel with an initial spectral efficiency
of 3\,b/s/Hz and up to four transmissions, modelling decoding as successful when the accumulated mutual
information (IR) or the combined SNR (Chase) supports the rate.}

\begin{inpractice}[HARQ in LTE and 5G]
LTE uses \emph{synchronous} HARQ in the uplink and asynchronous HARQ in the downlink, with eight parallel
stop-and-wait processes in FDD so that the channel stays busy while each process waits for its
acknowledgement roughly 8\,ms later. Each transport block is turbo coded at rate 1/3 into a circular buffer,
and each transmission reads a segment of the buffer starting at one of four \emph{redundancy versions},
conventionally sent in the order RV0, RV2, RV3, RV1: incremental redundancy with a single mother code.
5G NR keeps the idea with LDPC base graphs designed for it, allows up to 16 HARQ processes, flexible timing
and retransmission of only the failed \emph{code-block groups} of a large transport block. The details
are in Chapter~\ref{ch:cellular}.
\end{inpractice}

%======================================================================
\section{Choosing a classical code}
\label{sec:ch14:choosing}
%======================================================================

The codes of this chapter are not museum pieces. They survive wherever their particular strengths
matter more than the last decibel: tiny blocks, deterministic latency, simple hardware, burst errors,
erasures, or the need for a guaranteed rather than a statistical correction capability.
Table~\ref{tab:ch14:where} collects the main families and where an engineer will meet them.

\begin{table}[htbp]\centering\small
\caption{Where the classical codes live.}
\label{tab:ch14:where}
\begin{tabularx}{\linewidth}{@{}l l X@{}}
\toprule
Code & Strength & Typical systems\\
\midrule
Parity, CRC & cheap detection, burst detection & every packet network, storage, HARQ, polar-code list decoding\\
Hamming / SECDED & single-error correction in one gate delay & DRAM and cache ECC, on-die ECC, buses\\
Golay & triple-error correction in short words & HF ALE, space-probe imaging (Voyager)\\
BCH & $t$ random errors, simple hard decoder & DVB-S2/T2 outer code, flash (pre-LDPC), paging\\
Reed--Solomon & MDS, bursts, erasures & CD/DVD/Blu-ray, QR, DVB, ATSC, CCSDS, RAID-6, Ethernet KR4/KP4, OTN\\
Reed--Muller & fast soft ML decoding, short blocks & Mariner, LTE/NR short control reports\\
Convolutional + Viterbi & soft decisions, streaming, small latency & 802.11, DVB-S/T, CCSDS, LTE control, GNSS\\
Concatenated RS + conv. & very low error rates & deep space (legacy), DVB-S, DVB-T\\
Product codes & large distance from small codes & DVD, optical transport (staircase)\\
\bottomrule
\end{tabularx}
\end{table}

A few rules of thumb summarise decades of engineering experience. If a return channel exists and
latency is tolerable, detect with a CRC and retransmit, adding FEC only to bring the frame error rate
down to where retransmissions are rare. If errors come in bursts, use symbols (Reed--Solomon) or
interleave, and remember that interleaving costs delay. If soft information is available, use a
decoder that can exploit it, which in the classical toolbox means a convolutional code and the Viterbi
algorithm. If a hard guarantee on the number of correctable errors is needed, as in memory, use an
algebraic code. And if the last one or two decibels matter---as they do in every modern broadband
wireless and satellite system---move on to the codes of Chapter~\ref{ch:moderncodes}, keeping a CRC,
and often a short BCH or RS outer code, as insurance.

%======================================================================
\section{Summary}
%======================================================================

\begin{itemize}
\item Coding trades bandwidth, latency and complexity for power. Coding gain is measured at equal
energy per information bit and a target error rate; the asymptotic gain is $10\log_{10}(Rd)$ with soft
decisions and roughly 3\,dB less with hard decisions; at practical error rates soft decisions are worth
about 2\,dB.
\item A linear $(n,k,d)$ block code is defined by a generator matrix $\vect{G}$ or a parity-check matrix
$\vect{H}$. The syndrome $\vect{rH}^{\mathsf T}$ depends only on the error pattern; a code corrects $t$
errors and $e$ erasures when $2t+e<d$.
\item The Hamming, Singleton, Gilbert--Varshamov and Plotkin bounds frame what is possible. Hamming and
Golay codes are perfect; Reed--Solomon codes are MDS.
\item Extended Hamming (SECDED) codes protect computer memory; the (23,12) Golay code corrects three errors
in 23 bits.
\item Finite fields GF($2^m$) are built from a primitive polynomial; elements add by XOR and multiply by
adding logarithms. The choice of polynomial and bit conventions must match exactly between encoder and
decoder.
\item Cyclic codes are generated by a polynomial $g(x)$ dividing $x^n-1$ and encoded by an LFSR. A CRC with
$r$ check bits detects all bursts up to length $r$ and misses about $2^{-r}$ of random errors; its
performance at a given frame length is set by its Hamming distance there.
\item BCH codes place $2t$ consecutive roots in $g(x)$ to guarantee $d\ge2t+1$; they are decoded by
syndromes, the Berlekamp--Massey algorithm and a Chien search. Reed--Solomon codes do the same over
GF($2^m$) symbols and add Forney's formula for the error values; they correct bursts and erasures
superbly and are everywhere from CDs to Ethernet.
\item Interleaving spreads bursts across codewords at the cost of delay; concatenating a soft-decoded inner
convolutional code with an RS outer code reached $10^{-5}$ at about 2.3\,dB, the state of the art until
1993. Product and Reed--Muller codes build long codes from short ones.
\item Convolutional codes are finite-state machines described by octal generators, state diagrams and
trellises; the free distance and transfer function determine their performance, and the $K=7$ $(171,133)$
code with soft Viterbi decoding gains about 5\,dB at $10^{-5}$. Puncturing adjusts the rate and tail-biting
removes the termination overhead.
\item The Viterbi algorithm finds the maximum-likelihood path with $2^{K-1}$ add--compare--select operations
per bit; three-bit quantisation and a decision depth of about $5K$ lose almost nothing. Sequential
decoding handles long constraint lengths but is limited by the cutoff rate $R_0$.
\item ARQ uses a CRC and a return channel; selective repeat is the most efficient. HARQ with incremental
redundancy combines FEC and ARQ into a rate-adaptive scheme that tracks capacity without channel knowledge
at the transmitter.
\end{itemize}

\begin{labbox}[Block codes, CRCs, Reed--Solomon and Viterbi decoding]
Two notebooks accompany this chapter:
\begin{itemize}
\item \lab{moderncomms-labs/labs/lab08\_convolutional.py}
\item \lab{moderncomms-labs/labs/lab20\_block\_codes.py}
\end{itemize}
The first measures the Shannon limit and the constrained capacity of real constellations
(Chapter~\ref{ch:infotheory}), decodes Hamming (7,4) by syndromes, tests a 5G NR CRC-11 against single errors,
bursts and random corruption, simulates the $K=7$ $(133,171)$ code with hard and soft Viterbi decoding
(Figure~\ref{fig:ch14_hard_soft}), punctures it to rate 3/4 as in 802.11a, and shows how a block interleaver
rescues the decoder on a bursty Gilbert--Elliott channel. The second builds the algebraic codes: Hamming and
SECDED encoders with syndrome decoding and a miscorrection experiment, a CRC calculator for the standard
polynomials of Table~\ref{tab:ch14:crc} with all the reflection and initialisation options, a GF($2^m$)
playground that prints log/antilog tables, and Reed--Solomon encoding and decoding with random errors,
bursts, erasures and interleaving (Figures~\ref{fig:ch14_rs_performance} and~\ref{fig:ch14_rs_image}). Both
use \lab{commlib/coding.py} and the new \lab{commlib/gf.py}, whose Berlekamp--Massey, Chien and Forney
routines follow Section~\ref{sec:ch14:rs} line by line. The figure script
\lab{book/figscripts/ch14\_figs.py} contains a vectorised Viterbi decoder with truncated traceback and
quantisation options, and a transfer-function calculator for any feed-forward convolutional code.
\end{labbox}

\problems
\begin{enumerate}[label=\thechapter.\arabic*]
\item Compare the rate-1/3 repetition code with uncoded BPSK at equal $E_b/N_0$, using hard-decision
majority voting. Show that the repetition code is worse at every $E_b/N_0$. Does soft-decision combining
of the three copies change the conclusion?
\item For the Hamming (7,4) code of~\eqref{eq:ch14:h74}, list all 16 codewords and verify the weight
distribution $A_3=A_4=7$, $A_7=1$. Decode $\vect{r}=1101\,101$ and $\vect{r}=0110\,011$ by syndrome.
\item Show that the (23,12) Golay code meets the Hamming bound with equality. Is there a perfect binary
code with $n=90$, $k=78$, $t=2$? (Check the sphere-packing condition; then look up why no such code exists.)
\item A memory system uses the (72,64) SECDED code. The raw bit error probability per read is
$10^{-12}$, independent across bits. Estimate the probabilities per 64-bit word of (a) a corrected error,
(b) a detected uncorrectable error and (c) a silent miscorrection. Why are real miscorrection rates dominated
by multi-bit upsets rather than by this calculation?
\item In GF(16) with $p(x)=x^4+x+1$, compute $\alpha^7+\alpha^{11}$, $\alpha^7\alpha^{11}$,
$(\alpha^{7})^{-1}$ and $(\alpha^5+1)/(\alpha^3+\alpha)$. Verify that $x^4+x^3+x^2+x+1$ is irreducible but
not primitive.
\item Compute the CRC of the message \texttt{1011001110} with $g(x)=x^4+x^3+1$ by long division and with
the LFSR of Figure~\ref{fig:ch14_lfsr} (adapted to this generator). Which error patterns of weight two in
the 14-bit frame are undetected?
\item Show that if $g(x)=(x+1)p(x)$ with $p(x)$ primitive of degree $r-1$, the CRC detects every error
pattern of weight 1, 2 or 3 in frames shorter than $2^{r-1}-1$ bits. What is its Hamming distance at those
lengths?
\item Find the generator polynomial of the triple-error-correcting (15,5) BCH code and verify that its
roots include $\alpha,\dots,\alpha^6$. Decode the received word $r(x)=x^{12}+x^5+x$ assuming the
all-zero codeword was sent, using the Berlekamp--Massey algorithm.
\item For RS(7,3) over GF(8) as in the worked example, encode the message $(\alpha,0,1)$. Introduce one
error and two erasures and decode. What is the maximum number of erasures alone that the code can
correct?
\item Derive the transfer function~\eqref{eq:ch14:tf} of the $(7,5)$ code from the state diagram of
Figure~\ref{fig:ch14_state}, including the path-length variable $L$. How many error events of weight 7 are
there, and how long are they?
\item Show that the rate-1/2 code with generators $(3,5)_8$ is catastrophic by finding an input sequence of
infinite weight whose codeword has finite weight. Then use the Massey--Sain condition to check the codes of
Table~\ref{tab:ch14:conv}.
\item Using the union bound~\eqref{eq:ch14:ub} and the $B_d$ of Table~\ref{tab:ch14:conv}, find the
$E_b/N_0$ needed for $P_b=10^{-6}$ with the $K=3$, 5 and 7 codes. How much does each increase of two in
$K$ buy, and what does it cost in ACS operations per bit?
\item Stop-and-wait ARQ runs over a geostationary satellite link (round trip about 0.5\,s) at 2\,Mb/s with
1000-byte frames and a frame error rate of 1\%. Compute its efficiency, and the window size and efficiency of
go-back-$N$ and selective repeat.
\item \emph{(Simulation)} Using \texttt{commlib.gf}, measure the frame error rate of RS(255,223) decoding when
the symbol errors come from a Gilbert--Elliott channel, with interleaving depths of 1, 2, 4 and 8. Compare with
the formula for independent errors.
\item \emph{(Simulation)} Implement a Viterbi decoder for the $K=3$ $(7,5)$ code from scratch with Hamming
metrics, decode the worked example, and then add soft metrics and a decision depth. Plot the bit error rate
against $E_b/N_0$ and against decision depth, and compare with Figure~\ref{fig:ch14_viterbi_practical}.
\end{enumerate}

\furtherreading
\begin{itemize}
\item S.~Lin and D.~J.~Costello, Jr., \emph{Error Control Coding}, 2nd ed., Pearson Prentice Hall, 2004. The
standard engineering text: block, cyclic, BCH, RS, convolutional and trellis codes, and ARQ, with modern codes
added in this edition.
\item F.~J.~MacWilliams and N.~J.~A.~Sloane, \emph{The Theory of Error-Correcting Codes}, North-Holland, 1977.
The encyclopaedic mathematical reference; the source for weight enumerators, the MacWilliams identity and the
Golay and Reed--Muller codes.
\item R.~E.~Blahut, \emph{Algebraic Codes for Data Transmission}, Cambridge University Press, 2003. Finite
fields and algebraic decoding explained for engineers, with the Berlekamp--Massey and Euclidean algorithms in
detail.
\item T.~K.~Moon, \emph{Error Correction Coding: Mathematical Methods and Algorithms}, Wiley, 2005. A
programmer-friendly treatment with pseudocode for nearly every algorithm in this chapter.
\item S.~B.~Wicker and V.~K.~Bhargava (eds.), \emph{Reed--Solomon Codes and Their Applications}, IEEE Press,
1994. Chapters by the engineers who put RS codes into deep space, compact discs and storage.
\item A.~J.~Viterbi, ``Error bounds for convolutional codes and an asymptotically optimum decoding
algorithm,'' \emph{IEEE Transactions on Information Theory}, vol.~13, no.~2, 1967; and G.~D.~Forney, Jr.,
``The Viterbi algorithm,'' \emph{Proceedings of the IEEE}, vol.~61, no.~3, 1973, the tutorial that explained it
to everyone else.
\item R.~Johannesson and K.~Sh.~Zigangirov, \emph{Fundamentals of Convolutional Coding}, IEEE Press, 1999. The
definitive treatment of convolutional codes, distance properties and sequential decoding.
\item D.~J.~Costello, Jr.\ and G.~D.~Forney, Jr., ``Channel coding: the road to channel capacity,''
\emph{Proceedings of the IEEE}, vol.~95, no.~6, 2007. A superb historical survey from Hamming to LDPC codes,
including the deep-space milestones.
\item R.~W.~Hamming, ``Error detecting and error correcting codes,'' \emph{Bell System Technical Journal},
vol.~29, no.~2, 1950; I.~S.~Reed and G.~Solomon, ``Polynomial codes over certain finite fields,'' \emph{Journal
of the SIAM}, vol.~8, no.~2, 1960. Two short papers that are still a pleasure to read.
\item P.~Koopman, ``32-bit cyclic redundancy codes for Internet applications,'' \emph{Proc.\ International
Conference on Dependable Systems and Networks}, 2002; and R.~N.~Williams, ``A painless guide to CRC error
detection algorithms,'' 1993. How to choose a CRC, and how to implement it without tears.
\item CCSDS 131.0-B, \emph{TM Synchronization and Channel Coding} (Blue Book), current issue; ETSI EN~300~421
(DVB-S) and EN~302~307 (DVB-S2). The standards that define the RS, convolutional, concatenated and BCH codes
discussed here, with their exact polynomials and interleavers.
\end{itemize}
